% !TEX root = ../../main.tex
\setchapterimage{figure/chapter3封面图.jpg}
\setchapterintro{}
\chapter{线性映射，矩阵与行列式}
\section{向量空间上的线性映射}

\begin{definition}\phantomsection\label{p1:concept:c3-l7}
	对映射 $f:\F^n\to\F^m$，$\forall\alpha,\beta\in\F^n$，若 $f$ 满足
	\begin{enumerate}
		\item $f(\alpha+\beta)=f(\alpha)+f(\beta)$；
		\item $f(k\alpha)=k f(\alpha)$，$k\in\F$，
	\end{enumerate}
	就称 $f$ 为一个线性映射。用 $\operatorname{Hom}(\F^n,\F^m)$ 表示所有从 $\F^n$ 到 $\F^m$ 的线性映射构成的集合，并且
	\[
	\operatorname{End}(\F^n)=\operatorname{Hom}(\F^n,\F^n)
	\]
	
	零映射
	\[
	0:\alpha\in\F^n\longmapsto 0\in\F^m
	\]
	满足线性性，故 $0\in\operatorname{Hom}(\F^n,\F^m)$，因此 $\operatorname{Hom}(\F^n,\F^m)$ 不是空集。
	

\end{definition}
	在 $\F^n$ 中，取单位向量
\[
\varepsilon_1=\begin{pmatrix}1\\0\\\vdots\\0\end{pmatrix},
\quad \cdots \quad,
\varepsilon_n=\begin{pmatrix}0\\0\\\vdots\\1\end{pmatrix}.
\]
在 $\F^m$ 中，相应也有
\[
\varepsilon'_1=\begin{pmatrix}1\\0\\\vdots\\0\end{pmatrix},
\quad \cdots \quad,
\varepsilon'_m=\begin{pmatrix}0\\0\\\vdots\\1\end{pmatrix}.
\]
$\forall\alpha\in\F^n$，
\[
\alpha=a_1\varepsilon_1+\cdots+a_n\varepsilon_n
=\begin{pmatrix}a_1\\\vdots\\a_n\end{pmatrix}.
\]
由线性性，
\[
f(\alpha)=a_1f(\varepsilon_1)+\cdots+a_nf(\varepsilon_n).
\]
所以一个线性映射 $f$ 可以完全由它单位向量的像决定。

在 $\varepsilon'_1,\cdots,\varepsilon'_m$ 下，我们表示
\[
\left\{
\begin{aligned}
	f(\varepsilon_1)&=(a_{11},\cdots,a_{m1})
	=a_{11}\varepsilon'_1+\cdots+a_{m1}\varepsilon'_m,\\
	f(\varepsilon_2)&=(a_{12},\cdots,a_{m2})
	=a_{12}\varepsilon'_1+\cdots+a_{m2}\varepsilon'_m,\\
	&\vdots\\
	f(\varepsilon_n)&=(a_{1n},\cdots,a_{mn})
	=a_{1n}\varepsilon'_1+\cdots+a_{mn}\varepsilon'_m.
\end{aligned}
\right.
\]
把 $f(\varepsilon_j)$ 写成一个列向量
\[
\begin{pmatrix}a_{1j}\\\vdots\\a_{mj}\end{pmatrix}.
\]
将所有这样的向量放在一起，我们得到一个矩阵
\[
A=\bigl(f(\varepsilon_1)\ f(\varepsilon_2)\ \cdots\ f(\varepsilon_n)\bigr)
=\begin{pmatrix}
	a_{11}&a_{12}&\cdots&a_{1n}\\
	a_{21}&a_{22}&\cdots&a_{2n}\\
	\vdots&\vdots&\ddots&\vdots\\
	a_{m1}&a_{m2}&\cdots&a_{mn}
\end{pmatrix}.
\]
\phantomsection\label{p1:concept:c3-l77}我们称 $A$ 为 $f$ 的矩阵，或者说矩阵 $A$ 决定了一个线性映射 $f$，记为 $f_A$。
自然地，我们得到一个映射
\[
\rho:M_{m\times n}(\F)\longrightarrow \operatorname{Hom}(\F^n,\F^m),
\quad A\longmapsto f_A.
\]
\begin{proposition}
	上述的映射 $\rho$ 是单射且是满射。
\end{proposition}
\begin{proof}
	先证 $\rho$ 是单射，即证明 $\rho(A)=\rho(B)\Rightarrow A=B$。
	设 $f_A=f_B$，即 $\forall x\in\F^n$，都有
	\[
	Ax=Bx.
	\]
	取 $x=\varepsilon_i\in\F^n$，$i=1,\cdots,n$，那么
	\[
	A\varepsilon_i=B\varepsilon_i.
	\]
	即 $A$ 与 $B$ 的第 $i$ 列相同，故 $A=B$。
\end{proof}
关于$\rho$是满射的证明，我们需要一些准备工作
\begin{proposition}(线性映射的性质)
	设 $f\in\operatorname{Hom}(\F^n,\F^m)$，则
	\begin{enumerate}
		\item $f(0)=0\in\F^m$；
		\item $f\left(\sum_{i=1}^{s}k_i\alpha_i\right)=\sum_{i=1}^{s}k_i f(\alpha_i)$；
		若 $\alpha_1,\cdots,\alpha_s$ 线性相关，则 $f(\alpha_1),\cdots,f(\alpha_s)$ 亦线性相关；
		\item $f$ 保无关，当且仅当 $f$ 是单射。
	\end{enumerate}
\end{proposition}
\begin{proof}
	$1,2$ 显然。
	
	$3.$ 必要性：若 $\alpha_1\neq\alpha_2$，则 $\alpha_1-\alpha_2\neq0$，而 $\{\alpha_1-\alpha_2\}$ 线性无关。
	故 $\{f(\alpha_1-\alpha_2)\}$ 线性无关，也即
	\[
	f(\alpha_1-\alpha_2)\neq0.
	\]
	于是 $f(\alpha_1)\neq f(\alpha_2)$，故 $f$ 是单射。
	
	充分性：对无关组 $\alpha_1,\cdots,\alpha_n$，考虑
	\[
	k_1f(\alpha_1)+\cdots+k_nf(\alpha_n)=0.
	\]
	则
	\[
	f(k_1\alpha_1+
	\cdots+k_n\alpha_n)=0.
	\]
	由 $f$ 是单射，
	\[
	k_1\alpha_1+\cdots+k_n\alpha_n=0.
	\]
	故 $k_1=\cdots=k_n=0$，因此 $f(\alpha_1),\cdots,f(\alpha_n)$ 线性无关。
\end{proof}

\begin{lemma}
	$\varepsilon_1,\cdots,\varepsilon_n$ 是 $\F^n$ 的一组单位向量。
	\begin{enumerate}
		\item 若 $f,g\in\operatorname{Hom}(\F^n,\F^m)$，且
		\[
		f(\varepsilon_i)=g(\varepsilon_i),\quad i=1,\cdots,n,
		\]
		则 $f=g$；
		\item $\forall f\in\operatorname{Hom}(\F^n,\F^m)$，且
		\[
		f(\varepsilon_i)=\begin{pmatrix}a_{1i}\\\vdots\\a_{mi}\end{pmatrix},
		\]
		则
		\[
		f\left(\begin{pmatrix}x_1\\\vdots\\x_n\end{pmatrix}\right)
		=
		\begin{pmatrix}
			a_{11}x_1+\cdots+a_{1n}x_n\\
			\vdots\\
			a_{m1}x_1+
			\cdots+a_{mn}x_n
		\end{pmatrix}.
		\]
	\end{enumerate}
\end{lemma}
\begin{proof}
	$1.$ 对任意 $\alpha=\sum k_i\varepsilon_i$，有
	\[
	f(\alpha)=\sum k_i f(\varepsilon_i)
	=\sum k_i g(\varepsilon_i)
	=g\left(\sum k_i\varepsilon_i\right)=g(\alpha),
	\]
	故 $f=g$。
	
	$2.$ 因为
	\[
	\begin{pmatrix}x_1\\\vdots\\x_n\end{pmatrix}
	=x_1\varepsilon_1+\cdots+x_n\varepsilon_n,
	\]
	所以
	\[
	\begin{aligned}
		f\left(\begin{pmatrix}x_1\\\vdots\\x_n\end{pmatrix}\right)
		&=x_1f(\varepsilon_1)+\cdots+x_nf(\varepsilon_n)\\
		&=x_1\begin{pmatrix}a_{11}\\\vdots\\a_{m1}\end{pmatrix}
		+\cdots+
		x_n\begin{pmatrix}a_{1n}\\\vdots\\a_{mn}\end{pmatrix}\\
		&=\begin{pmatrix}
			a_{11}x_1+\cdots+a_{1n}x_n\\
			\vdots\\
			a_{m1}x_1+\cdots+a_{mn}x_n
		\end{pmatrix}.
	\end{aligned}
	\]
\end{proof}

下面，证明 $\rho$ 是一个满射。
\begin{proof}[证明]
	$\rho$ 是满射，也即 $\forall T\in\operatorname{Hom}(\F^n,\F^m)$，存在 $A\in M_{m\times n}(\F)$，使得
	\[
	f_A=T.
	\]
	$\forall T\in\operatorname{Hom}(\F^n,\F^m)$，取 $\varepsilon_i\in\F^n$，$i=1,\cdots,n$，计算 $T\varepsilon_i$，构造矩阵
	\[
	A=\bigl(T\varepsilon_1\ \cdots\ T\varepsilon_n\bigr).
	\]
	下证 $f_A=T$，只需 $\forall i$，$f_A(\varepsilon_i)=T\varepsilon_i$。
	事实上，
	\[
	f_A(\varepsilon_i)=A\varepsilon_i
	=\bigl(T\varepsilon_1\ \cdots\ T\varepsilon_n\bigr)\varepsilon_i
	=T\varepsilon_i.
	\]
\end{proof}

\begin{definition}\phantomsection\label{p1:concept:c3-l209}
	$f\in\operatorname{Hom}(\F^n,\F^m)$，定义 $f$ 的核 $\ker f$ 为
	\[
	\ker f:=f^{-1}(0)=\{\alpha\in\F^n:f(\alpha)=0\}.
	\]
	显然，$0\in\ker f$，故 $\varnothing\neq\ker f$。
\end{definition}

\begin{lemma}\label{3.1.6}
	$f\in\operatorname{Hom}(\F^n,\F^m)$，则 $f$ 是单射，当且仅当 $\ker f=\{0\}$。
\end{lemma}
\begin{proof}
	必要性：若 $\alpha\in\F^n$ 且 $f(\alpha)=f(0)=0$，则由 $f$ 是单射，$\alpha=0$，故 $\ker f=\{0\}$。
	
	充分性：若 $f(\alpha_1)=f(\alpha_2)$ 且 $\alpha_1\neq\alpha_2$，则
	\[
	f(\alpha_1-\alpha_2)=0.
	\]
	由 $\ker f=\{0\}$，$\alpha_1-\alpha_2=0$，故 $\alpha_1=\alpha_2$，矛盾。因此 $f$ 是单射。
\end{proof}

考虑由 $f$ 生成的矩阵 $A$，不难得到以下结论。
\begin{proposition}\label{3.1.7}
	$f\in\operatorname{Hom}(\F^n,\F^m)$，且 $A\in M_{m\times n}(\F)$ 是 $f$ 生成的矩阵。
	那么，对 $\alpha\in\F^n$，$\alpha\in\ker f$ 当且仅当 $\alpha$ 是以 $A$ 为系数矩阵的齐次线性方程组的解，也即
	\[
	Ax=0\text{ 的解空间}=\ker f.
	\]
\end{proposition}
\begin{proof}
	显然。
\end{proof}

利用 \ref{2.4.12}、\ref{3.1.6} 与 \ref{3.1.7}，我们得到以下结论。
\begin{corollary}
	$f$ 是单射，当且仅当 $r(A)=n$。
\end{corollary}
\phantomsection\label{p1:concept:c3-l246}我们称 $r(A)=n$ 的情况为列满秩的。

\begin{definition}\phantomsection\label{p1:concept:c3-l248}
	对 $f\in\operatorname{Hom}(\F^n,\F^m)$，定义 $f$ 的像 $\operatorname{im}(f)$ 为
	\[
	\operatorname{im}(f):=\{\beta\in\F^m:\exists\alpha\in\F^n\text{ 使得 } f(\alpha)=\beta\}.
	\]
	考虑 $f$ 对应的 $A\in M_{m\times n}(\F)$，若 $\beta\in\operatorname{im}(f_A)$，当且仅当 $\exists x\in\F^n$ 使得
	\[
	Ax=\beta,
	\]
	也即 $\beta$ 可由 $A$ 的列向量线性表示。
\end{definition}

\begin{proposition}
	$f_A$ 是满射（或 $\operatorname{im}(f_A)=\F^m$），当且仅当 $r(A)=m$。
\end{proposition}
\phantomsection\label{p1:concept:c3-l263}我们称 $r(A)=m$ 的情况为行满秩的。
\begin{proof}
	必要性：
	\[
	\operatorname{im}(f_A)=\F^m
	\Longleftrightarrow
	\forall\beta\in\F^m\text{ 可由 }\alpha_1,\cdots,\alpha_n\text{ 表出}.
	\]
	取 $\F^m$ 的单位向量 $\varepsilon_1,\cdots,\varepsilon_m$，那么 $\varepsilon_1,\cdots,\varepsilon_m$ 可由 $\alpha_1,\cdots,\alpha_n$ 表出。
	由基本引理，
	\[
	r(\varepsilon_1,\cdots,\varepsilon_m)=m
	\leq r(\alpha_1,\cdots,\alpha_n)\leq m.
	\]
	故
	\[
	r(\alpha_1,\cdots,\alpha_n)=m.
	\]
	充分性：已知 $r(A)=m$，证明 $\beta\in\F^m$ 都能找到
	\[
	x_1\alpha_1+\cdots+x_n\alpha_n=\beta
	\]
	的解 $x_1,\cdots,x_n$。那么
	\[
	r(\alpha_1,\cdots,\alpha_n)=r(\alpha_1,\cdots,\alpha_n,\beta)=m,
	\]
	故有解。
	

\end{proof}
	因此，$f_A$ 既是单射又是满射，当且仅当
\[
r(A)=m=n.
\]

\section{线性映射的运算与矩阵的运算}

\begin{definition}\phantomsection\label{p1:concept:c3-l300}
	$\operatorname{Hom}(\F^n,\F^m)$ 上的加法与数乘定义如下。
	对 $f,g\in\operatorname{Hom}(\F^n,\F^m)$，$\alpha\in\F^n$，定义
	\[
	f+g:=\alpha\longmapsto f(\alpha)+g(\alpha),
	\]
	也即
	\[
	(f+g)(\alpha)=f(\alpha)+g(\alpha).
	\]
	显然，也有
	\[
	f+g\in\operatorname{Hom}(\F^n,\F^m).
	\]
	对 $k\in\F$，定义
	\[
	kf:=\alpha\longmapsto kf(\alpha),
	\]
	则
	\[
	kf\in\operatorname{Hom}(\F^n,\F^m).
	\]
	
\end{definition}
因此 $\operatorname{Hom}(\F^n,\F^m)$ 在上述的加法与数乘上是一个线性空间(Chapter 4)。

下面考虑
映射的复合。
\begin{definition}\phantomsection\label{p1:concept:c3-l328}
	设
	\[
	g\in\operatorname{Hom}(\F^n,\F^m),\qquad
	f\in\operatorname{Hom}(\F^m,\F^\ell),\qquad
	\alpha\in\F^n.
	\]
	定义
	\[
	fg:=\alpha\longmapsto f(g(\alpha)).
	\]
	这里 $\alpha\in\F^n$，$g(\alpha)\in\F^m$，$f(g(\alpha))\in\F^\ell$。
	显然，也有
	\[
	fg\in\operatorname{Hom}(\F^n,\F^\ell).
	\]
	
\end{definition}
考虑 $fg,g,f$ 的核与像，有
\[
\ker g\subset \ker(fg),
\qquad
\operatorname{im}(fg)\subset \operatorname{im}f.
\]
也就有
\[
fg\text{ 单射}\Rightarrow g\text{ 单射},
\qquad
fg\text{ 满射}\Rightarrow f\text{ 满射}.
\]
当 $n=m=\ell$ 时，$f,g\in\operatorname{End}(\F^n)$，$fg$ 称为 $\F^n$ 上的乘法。
\phantomsection\label{p1:concept:c3-l359}$\operatorname{End}(\F^n)$ 是一个环，同称为 $\F^n$ 的自同态环。
加上数乘，$\operatorname{End}(\F^n)$ 是一个结合代数，$n\geq 2$ 时通常是非交换的。

对 $f\in\operatorname{Hom}(\F^n,\F^m)$，$g\in\operatorname{Hom}(\F^m,\F^\ell)$，当 $n,m,\ell$ 中至少有一个大于 $1$ 时，$fg$ 与 $gf$ 一般也不满足交换律。
\begin{definition}\phantomsection\label{p1:concept:c3-l363}
	对 $f\in\operatorname{End}(\F^n)$，定义 $f$ 的幂与多项式。
	\[
	f^2=ff,
	\qquad
	f^3=fff,
	\qquad \cdots
	\]
	就有链
	\[
	\ker f\subset \ker f^2\subset \ker f^3\subset\cdots,
	\]
	\[
	\operatorname{im}(f)\supset \operatorname{im}(f^2)\supset \operatorname{im}(f^3)\supset\cdots.
	\]
	规定
	\[
	f^0=\operatorname{id}_{\F^n}:\alpha\longmapsto \alpha.
	\]
	对
	\[
	h(x)=a_0+a_1x+\cdots+a_mx^m\in\F[x],
	\]
	定义
	\[
	h(f):=a_0\operatorname{id}+a_1f+\cdots+a_mf^m\in\operatorname{End}(\F^n).
	\]
	称 $h(f)$ 为 $f$ 的一个多项式。$\forall h\in\F[x]$，$h(f)$ 与 $f$ 是可交换的。
\end{definition}

对矩阵的运算，我们采取映射的视角。

\begin{definition}[矩阵加法]\phantomsection\label{p1:concept:c3-l395}
	设 $A,B\in M_{m\times n}(\F)$，$f,g$ 分别是 $A,B$ 对应的映射，
	$f,g\in\operatorname{Hom}(\F^n,\F^m)$。取 $\F^n$ 的单位向量
	$\varepsilon_1,\cdots,\varepsilon_n$，那么
	\[
	A=\bigl(f(\varepsilon_1)\ \cdots\ f(\varepsilon_n)\bigr),
	\qquad
	B=\bigl(g(\varepsilon_1)\ \cdots\ g(\varepsilon_n)\bigr).
	\]
	定义
	\[
	A+B:=\bigl(f(\varepsilon_1)+g(\varepsilon_1)\ \cdots\
	f(\varepsilon_n)+g(\varepsilon_n)\bigr).
	\]
	若
	\[
	f(\varepsilon_i)=
	\begin{pmatrix}
		A(1,i)\\ \vdots\\ A(m,i)
	\end{pmatrix},
	\qquad
	g(\varepsilon_i)=
	\begin{pmatrix}
		B(1,i)\\ \vdots\\ B(m,i)
	\end{pmatrix},
	\]
	那么 $A+B$ 的 $(i,j)$ 元定义为
	\[
	(A+B)(i,j)=A(i,j)+B(i,j).
	\]
	这由 $\F$ 上加法的性质，不难得出
	$M_{m\times n}(\F)$ 对加法构成一个交换群，也即满足交换律、结合律，存在单位元与逆元。
	

\end{definition}
	若使用矩阵的行向量形式与列向量形式，$A+B$ 可以写成
\[
(\alpha_1\ \cdots\ \alpha_n)+(\gamma_1\ \cdots\ \gamma_n)
=(\alpha_1+\gamma_1\ \cdots\ \alpha_n+\gamma_n),
\]
或
\[
\begin{pmatrix}
	\beta_1\\ \vdots\\ \beta_m
\end{pmatrix}
+
\begin{pmatrix}
	\delta_1\\ \vdots\\ \delta_m
\end{pmatrix}
=
\begin{pmatrix}
	\beta_1+\delta_1\\ \vdots\\ \beta_m+\delta_m
\end{pmatrix}.
\]
\begin{proposition}
	若 $A,B\in M_{m\times n}(\F)$，那么
	\[
	r(A+B)\leq r(A)+r(B).
	\]
\end{proposition}

\begin{proof}
	记
	\[
	A=(\alpha_1\ \cdots\ \alpha_n),\qquad
	B=(\beta_1\ \cdots\ \beta_n),
	\]
	那么
	\[
	A+B=(\alpha_1+\beta_1\ \cdots\ \alpha_n+\beta_n).
	\]
	考虑向量组
	\[
	\begin{array}{ll}
		(\mathrm{I})&\alpha_1,\cdots,\alpha_n,\\
		(\mathrm{II})&\beta_1,\cdots,\beta_n,\\
		(\mathrm{III})&\alpha_1+\beta_1,\cdots,\alpha_n+\beta_n,\\
		(\mathrm{IV})&\alpha_1,\cdots,\alpha_n,\beta_1,\cdots,\beta_n.
	\end{array}
	\]
	$(\mathrm{III})$ 可由 $(\mathrm{IV})$ 表出，故
	\[
	r(A+B)
	\leq r(\alpha_1,\cdots,\alpha_n,\beta_1,\cdots,\beta_n)
	\leq r(A)+r(B).
	\]
\end{proof}

\begin{definition}[矩阵的数乘]\phantomsection\label{p1:concept:c3-l483}
	设 $A\in M_{m\times n}(\F)$，$k\in\F$，且 $A$ 对应的映射为 $f$。
	若
	\[
	A=\bigl(f(\varepsilon_1)\ \cdots\ f(\varepsilon_n)\bigr),
	\]
	则定义
	\[
	kA:=\bigl(kf(\varepsilon_1)\ \cdots\ kf(\varepsilon_n)\bigr).
	\]
	结合加法，矩阵的数乘满足
	\begin{enumerate}
		\item $1A=A$；
		\item 结合律：$k(\ell A)=(k\ell)A$；
		\item $(k+\ell)A=kA+\ell A$；
		\item $k(A+B)=kA+kB$。
	\end{enumerate}
\end{definition}

对 $B\in M_{m\times n}(\F)$，$A\in M_{\ell\times m}(\F)$，我们定义矩阵的乘法。

\begin{definition}[矩阵的乘法]\phantomsection\label{p1:concept:c3-l504}
	设 $B$ 对应于 $g$，$A$ 对应于 $f$，即
	\[
	\F^n\xrightarrow{\ g\leftrightarrow B\ }\F^m
	\xrightarrow{\ f\leftrightarrow A\ }\F^\ell,
	\qquad
	\F^n\xrightarrow{\ fg\ }\F^\ell.
	\]
	取单位向量
	\[
	\varepsilon_1,\cdots,\varepsilon_n\in\F^n,
	\qquad
	\eta_1,\cdots,\eta_m\in\F^m.
	\]
	则
	\[
	\begin{aligned}
		fg(\varepsilon_j)
		&=f(g(\varepsilon_j))
		=f\left(
		\begin{pmatrix}
			B(1,j)\\ \vdots\\ B(m,j)
		\end{pmatrix}
		\right)\\
		&=f\bigl(B(1,j)\eta_1+\cdots+B(m,j)\eta_m\bigr)\\
		&=B(1,j)f(\eta_1)+\cdots+B(m,j)f(\eta_m)\\
		&=B(1,j)
		\begin{pmatrix}
			A(1,1)\\ \vdots\\ A(\ell,1)
		\end{pmatrix}
		+\cdots+
		B(m,j)
		\begin{pmatrix}
			A(1,m)\\ \vdots\\ A(\ell,m)
		\end{pmatrix}\\
		&=
		\begin{pmatrix}
			\sum_{k=1}^{m}A(1,k)B(k,j)\\
			\vdots\\
			\sum_{k=1}^{m}A(\ell,k)B(k,j)
		\end{pmatrix}.
	\end{aligned}
	\]
	于是
	\[
	AB=
	\begin{pmatrix}
		\sum_{k=1}^{m}A(1,k)B(k,1)&\cdots&\sum_{k=1}^{m}A(1,k)B(k,n)\\
		\vdots&\ddots&\vdots\\
		\sum_{k=1}^{m}A(\ell,k)B(k,1)&\cdots&\sum_{k=1}^{m}A(\ell,k)B(k,n)
	\end{pmatrix}
	\in M_{\ell\times n}(\F).
	\]

\end{definition}
	矩阵乘法满足
\begin{enumerate}
	\item 结合律：$(AB)C=A(BC)$；
	\item 左右分配律：$A(B+C)=AB+AC$，$(B+C)A=BA+CA$；
	\item 与数乘的结合律：$kAB=A(kB)$；
	\item 在 $M_{n\times n}(\F)$ 中，存在单位元
	\[
	E_n=\begin{pmatrix}
		1& &0\\
		&\ddots& \\
		0& &1
	\end{pmatrix};
	\]
	\item 无交换律：$AB$ 一般不等于 $BA$，例如
	\[
	\begin{pmatrix}0&1\\0&0\end{pmatrix}
	\begin{pmatrix}1&0\\0&0\end{pmatrix}
	=
	\begin{pmatrix}0&0\\0&0\end{pmatrix},
	\qquad
	\begin{pmatrix}1&0\\0Ħ0\end{pmatrix}
	\begin{pmatrix}0&1\\0&0\end{pmatrix}
	=
	\begin{pmatrix}0&1\\0&0\end{pmatrix}.
	\]
	若 $AB=BA$，则称 $A,B$ 可交换；
	\item 无消去律，也即存在 $A,B\neq0$，使 $AB=0$，例如
	\[
	\begin{pmatrix}0&1\\0&0\end{pmatrix}
	\begin{pmatrix}0&1\\0&0\end{pmatrix}
	=
	\begin{pmatrix}0&0\\0&0\end{pmatrix}.
	\]
\end{enumerate}
采取二元运算的视角，$AB=C$ 可以理解成
\[
M_{\ell\times m}(\F)\times M_{m\times n}(\F)
\longrightarrow M_{\ell\times n}(\F).
\]
对于 $n$ 阶方阵构成的集合 $M_n(\F)$，对上述的加法、乘法构成一个非交换环；加上数乘，构成一个非交换代数。

采取行向量或列向量形式，我们不难得到以下命题。

\begin{proposition}
	若 $A\in M_{m\times n}(\F)$，$B\in M_{n\times \ell}(\F)$，那么
	\[
	r(AB)\leq \min\{r(A),r(B)\}.
	\]
\end{proposition}

\begin{proof}
	先设
	\[
	A=\begin{pmatrix}
		a_{11}&\cdots&a_{1n}\\
		\vdots&\ddots&\vdots\\
		a_{m1}&\cdots&a_{mn}
	\end{pmatrix},
	\qquad
	B=\begin{pmatrix}
		\beta_1\\ \vdots\\ \beta_n
	\end{pmatrix}.
	\]
	则
	\[
	AB=\begin{pmatrix}
		a_{11}\beta_1+\cdots+a_{1n}\beta_n\\
		\vdots\\
		a_{m1}\beta_1+\cdots+a_{mn}\beta_n
	\end{pmatrix}.
	\]
	于是
	\[
	\begin{aligned}
		r(AB)
		&=r(a_{11}\beta_1+
		\cdots+a_{1n}\beta_n,\cdots,
		a_{m1}\beta_1+
		\cdots+a_{mn}\beta_n)\\
		&\leq r(\beta_1,\cdots,\beta_n)=r(B).
	\end{aligned}
	\]
	再取 $A$ 的列向量形式即可得
	\[
	r(AB)\leq r(A).
	\]
	故
	\[
	r(AB)\leq \min\{r(A),r(B)\}.
	\]
\end{proof}

下面，我们给出 Kronecker 记号，并以此定义基本矩阵。
\phantomsection\label{p1:concept:c3-l652}Kronecker 记号是由两个元决定的 $\delta_{ij}$，定义为
\[
\delta_{ij}=\begin{cases}
	1,&i=j,\\
	0,&i\neq j.
\end{cases}
\]

\begin{definition}[基本矩阵]\phantomsection\label{p1:concept:c3-l660}
	在 $M_{m\times n}(\F)$ 上，基本矩阵 $E_{ij}$ 的 $(k,l)$ 元定义为
	\[
	E_{ij}(k,l)=\delta_{ik}\delta_{jl}.
	\]
	也即 $E_{ij}$ 是只有 $(i,j)$ 元为 $1$，其余位置都为 $0$ 的矩阵。
	

\end{definition}
	那么 $\forall A\in M_{m\times n}(\F)$，其 $(i,j)$ 元为 $A(i,j)$，就有
\[
A=\sum_{\substack{1\leq i\leq m\\1\leq j\leq n}} A(i,j)E_{ij}.\tag{1}
\]
也就是说，$A$ 可由 $E_{ij}$（$1\leq i\leq m,1\leq j\leq n$）线性表出。
可以看到，$E_{ij}$ 构成了一个矩阵意义下的极大无关组。
那么 $M_{m\times n}(\F)$ 中任意 $mn+1$ 个矩阵必然相关。

利用基本矩阵，我们可以得到另一种矩阵乘法的定义方式：
\[
E_{ij}E_{kl}=\delta_{jk}E_{il}.
\]
其中 $E_{ij}\in M_{s\times m}(\F)$，$E_{kl}\in M_{m\times n}(\F)$。
那么 $\forall A\in M_{s\times m}(\F)$，$B\in M_{m\times n}(\F)$，找出它们如 $(1)$ 的形式，再做乘法的分配即可定义 $AB$。
\begin{definition}(矩阵的转置)\phantomsection\label{p1:concept:c3-l683}
	设 $A\in M_{m\times n}(\F)$，且
	\[
	A=(A(i,j)).
	\]
	定义 $A$ 的转置 $A'$（或 $A^T$）为
	\[
	A'\in M_{n\times m}(\F),
	\qquad
	A'=(A(j,i)).
	\]
	特别地，在 $\mathbb C$ 上，我们有复共轭转置
	\[
	A^*=(\overline{A(j,i)}).
	\]
\end{definition}

\begin{proposition}[转置的性质]
	转置满足
	\begin{enumerate}
		\item $A'+B'=(A+B)'$；
		\item $kA'=(kA)'$；
		\item $(AB)'=B'A'$；
		\item $(A')'=A$。
	\end{enumerate}
\end{proposition}

\begin{proof}
	$1,2,4$ 显然。对 $3$，考虑基本矩阵即可，证明略。
\end{proof}


\section{方阵}

\begin{definition}\phantomsection\label{p1:concept:c3-l717}
	设 $A\in M_n(\F)$，那么 $A^k\in M_n(\F)$，$k\in\mathbb Z_+$。
	规定
	\[
	A^0=E_n=
	\begin{pmatrix}
		1&&0\\
		&\ddots&\\
		0&&1
	\end{pmatrix}.
	\]
	若存在 $k\in\mathbb Z_+$，使得
	\[
	A^k=0,
	\]
	则称 $A$ 是幂零的。
	
\end{definition}
类似于数的幂，矩阵的幂同样有
\[
(A^k)^\ell=A^{k\ell},
\qquad
A^kA^\ell=A^{k+\ell}.
\]
但是
\[
(AB)^k\neq A^kB^k,
\]
因为无交换律。
\begin{proposition}
	$\forall f,g\in\F[x]$ 且 $f,g\neq0$，都有
	\[
	f(A)g(A)=g(A)f(A).
	\]
\end{proposition}

\begin{proof}
	由多项式乘法的交换律与矩阵幂的交换关系得到。
\end{proof}

\begin{definition}\phantomsection\label{p1:concept:c3-l757}
	若 $f\neq0$ 且 $f(A)=0$，则称 $f$ 是 $A$ 的零化多项式。
	若 $f$ 是 $A$ 的零化多项式，则 $\forall k\in\F$，$kf$ 亦是 $A$ 的零化多项式。
\end{definition}

\begin{proposition}
	$\forall A\in M_n(\F)$，存在 $f\in\F[x]$，使得
	\[
	f(A)=0.
	\]
\end{proposition}

\begin{proof}
	考虑
	\[
	A^0,A^1,\cdots,A^{n^2}
	\]
	共 $n^2+1$ 个矩阵，必线性相关。故存在
	$k_0,\cdots,k_{n^2}$ 不全为 $0$，使得
	\[
	k_0A^0+k_1A^1+\cdots+k_{n^2}A^{n^2}=0.
	\]
	令
	\[
	f(x)=k_0+k_1x+\cdots+k_{n^2}x^{n^2}\in\F[x],
	\]
	且 $f\neq0$，故 $f(x)$ 为 $A$ 的零化多项式。
\end{proof}

\begin{proposition}
	\phantomsection\label{p1:concept:c3-l787}$\forall A\in M_n(\F)$，$A$ 的次数最低的且首一的零化多项式存在且唯一，称为 $A$ 的极小多项式。
\end{proposition}

\begin{proof}
	取 $m_1,m_2\in\F[x]$ 是 $A$ 的次数最低的且首一的零化多项式。
	做带余除法
	\[
	m_1=m_2q+r,
	\]
	于是
	\[
	0=m_1(A)=m_2(A)q(A)+r(A).
	\]
	若 $r(x)\neq0$，则由 $r(A)=0$ 且 $\deg r<\deg m_2$，矛盾。
	故 $r=0$，$m_2\mid m_1$。
	同理 $m_1\mid m_2$，故 $m_1=m_2$。
\end{proof}

极小多项式至少一次，$A$ 的任意零化多项式可被极小多项式整除。

显然地，$n$ 阶数量阵 $aE_n$ 与任意 $n$ 阶方阵可交换。
那么除了 $aE_n$ 外，其它的矩阵是否满足上述条件？Schur 引理告诉我们只有 $aE_n$ 与任意 $A$ 可交换。

\begin{lemma}(Schur 引理)
	若 $n$ 阶方阵 $A$ 与任意 $n$ 阶方阵可交换，则 $A$ 本身是数量阵。
\end{lemma}

\begin{proof}
	对 $A$，$\forall B\in M_n(\F)$，有
	\[
	AB=BA.
	\]
	取 $B=E_{ij}$，$1\leq i,j\leq n$ 即可。
	事实上，比较 $AE_{ij}=E_{ij}A$ 的各个元素，可得 $A$ 的非对角元全为 $0$，并且所有对角元相等，故
	\[
	A=aE_n.
	\]
\end{proof}

\begin{definition}\phantomsection\label{p1:concept:c3-l826}
	$\forall A\in M_n(\F)$，$A=(a_{ij})$，定义 $A$ 的迹 $\operatorname{tr}(A)$ 为
	\[
	\operatorname{tr}(A)=\sum_i a_{ii}.
	\]
	
\end{definition}
$\operatorname{tr}(A)$ 是一个映射
\[
\operatorname{tr}:M_n(\F)\longrightarrow\F.
\]
显然也有
\begin{enumerate}
	\item $\operatorname{tr}(A+B)=\operatorname{tr}(A)+\operatorname{tr}(B)$；
	\item $\operatorname{tr}(kA)=k\operatorname{tr}(A)$；
	\item $\operatorname{tr}(A)=\operatorname{tr}(A')$；
	\item $\operatorname{tr}(AB)=\operatorname{tr}(BA)$。
\end{enumerate}
进一步，
\[
\operatorname{tr}(ABC)=\operatorname{tr}(CAB)=\operatorname{tr}(BCA),
\]
但一般
\[
\operatorname{tr}(AB)\neq\operatorname{tr}(A)\operatorname{tr}(B).
\]
从 $4$ 来看，当 $A,B\in M_n(\F)$ 时，不能有
\[
AB-BA=E_n.
\]
\begin{definition}\phantomsection\label{p1:concept:c3-l856}
	设 $A\in M_n(\F)$。
	若
	\[
	A'=A,
	\]
	则称 $A$ 是对称阵；若
	\[
	A'=-A,
	\]
	则称 $A$ 是反对称阵。
	对称阵 $A$ 是 $E_{ij}+E_{ji}$ 与 $E_{ii}$ 的线性组合；反对称阵 $A$ 是 $E_{ij}-E_{ji}$ 的线性组合。
	显然，$\forall A\in M_n(\F)$，存在 $S,T$，$S'=S$，$T'=-T$，使得
	\[
	A=S+T.
	\]
	令
	\[
	S=\frac{A+A'}{2},
	\qquad
	T=\frac{A-A'}{2}
	\]
	即可。
\end{definition}

\begin{definition}(三角阵)\phantomsection\label{p1:concept:c3-l881}
	设 $A\in M_n(\F)$，若 $\forall i,j$，当 $i>j$ 时都有 $A(i,j)=0$，则称 $A$ 是一个上三角阵，即形如
	\[
	\begin{pmatrix}
		*&\cdots&*&\cdots&*\\
		0&*&\cdots& &\vdots\\
		\vdots&\ddots&\ddots&\ddots&\vdots\\
		0&\cdots&0&*&*\\
		0&\cdots&\cdots&0&*
	\end{pmatrix}.
	\]
	类似地，我们定义出下三角阵。
\end{definition}

$\F$ 上所有 $n$ 阶上三角阵关于加法、数乘、乘法封闭。
对加法与数乘显然，对乘法证明如下。
对 $i>j$，有
\[
\begin{aligned}
	(AB)(i,j)
	&=\sum_{k=1}^{n}A(i,k)B(k,j)\\
	&=\sum_{k<i}\underbrace{A(i,k)}_{0}B(k,j)
	+\sum_{k\geq i}A(i,k)\underbrace{B(k,j)}_{0}\\
	&=0.
\end{aligned}
\]

\begin{definition}\phantomsection\label{p1:concept:c3-l908}
	称非对角元都是 $0$ 的矩阵为对角阵，记作
	\[
	\operatorname{diag}[a_1,\cdots,a_n]=
	\begin{pmatrix}
		a_1&&0\\
		&\ddots&\\
		0&&a_n
	\end{pmatrix}.
	\]
\end{definition}

\begin{definition}(初等阵)\label{3.3.11}\phantomsection\label{p1:concept:c3-l920}
	由单位阵经过一次初等变换得到的矩阵称为一个初等阵。
	三类初等变换分别对应一类初等阵：
	\[
	P(i,j),\qquad P(a(i)+(j)),\qquad P(c(i)).
	\]

\end{definition}
	显然，有
\[
P(i,j)^2=E_n,
\]
\[
P(a(i)+(j))P(-a(i)+(j))=E_n,
\]
\[
P(c(i))P(c^{-1}(i))=E_n.
\]
\begin{theorem}
	$\forall M\in M_{m\times n}(\F)$，有
	\begin{enumerate}
		\item 对 $M$ 做一次初等行变换得到的矩阵，等于用相应的初等阵左乘 $M$ 得到的矩阵；
		\item 对 $M$ 做一次初等列变换得到的矩阵，等于用相应的初等阵右乘 $M$ 得到的矩阵。
	\end{enumerate}
\end{theorem}

\begin{proof}
	我们只对 $P(i,j)$ 的情况证明。
	对行变换，设
	\[
	A=
	\begin{pmatrix}
		\alpha_1\\
		\vdots\\
		\alpha_i\\
		\vdots\\
		\alpha_j\\
		\vdots\\
		\alpha_m
	\end{pmatrix}.
	\]
	则
	\[
	P(i,j)A=
	\begin{pmatrix}
		\alpha_1\\
		\vdots\\
		\alpha_j\\
		\vdots\\
		\alpha_i\\
		\vdots\\
		\alpha_m
	\end{pmatrix}.
	\]
	对列变换，设
	\[
	P(i,j)=
	\begin{pmatrix}
		\varepsilon_1\\
		\vdots\\
		\varepsilon_j\\
		\vdots\\
		\varepsilon_i\\
		\vdots\\
		\varepsilon_n
	\end{pmatrix},
	\qquad
	A=(\alpha_1\ \cdots\ \alpha_i\ \cdots\ \alpha_j\ \cdots\ \alpha_n).
	\]
	则
	\[
	AP(i,j)=(\alpha_1\ \cdots\ \alpha_j\ \cdots\ \alpha_i\ \cdots\ \alpha_n).
	\]
	其它两类初等变换同理。
\end{proof}

根据 Gauss 消元法以及上述定理，我们得到以下推论。
\begin{corollary}
	设 $M,N\in M_{m\times n}(\F)$。
	$M$ 与 $N$ 相抵，当且仅当存在 $m$ 阶初等阵 $P_1,\cdots,P_s$ 与 $n$ 阶初等阵 $Q_1,\cdots,Q_t$，使得
	\[
	M=P_1\cdots P_sNQ_1\cdots Q_t.
	\]
	特别地，$\forall M\in M_{m\times n}(\F)$，且 $r(M)=r$，存在 $m$ 阶初等阵 $P_1,\cdots,P_s$ 与 $n$ 阶初等阵 $Q_1,\cdots,Q_t$，使得
	\[
	M=P_1\cdots P_s
	\begin{pmatrix}
		E_r&0\\
		0&0
	\end{pmatrix}
	Q_1\cdots Q_t.
	\]
\end{corollary}

\section{矩阵的逆}

\begin{definition}\label{3.4.1}\phantomsection\label{p1:concept:c3-l1016}
	对 $A\in M_n(\F)$，若 $\exists B\in M_n(\F)$，使得
	\[
	AB=BA=E_n,
	\]
	则称 $A$ 是可逆的。
	
	对可逆矩阵 $A$，满足 $AB=BA=E_n$ 的矩阵是唯一的。证明使用同一法：设
	\[
	AB_1=B_1A=E_n=AB_2=B_2A,
	\]
	则
	\[
	B_1=B_1E_n=B_1AB_2=E_nB_2=B_2.
	\]
	这样，记 $B=A^{-1}$，称为 $A$ 的逆。
\end{definition}

对于逆矩阵，我们给出以下这些经典的例子。
\begin{example}\label{3.4.2}
	\begin{enumerate}
		\item $A$ 可逆，则
		\[
		f_A:x\longmapsto Ax
		\]
		有逆映射，显然其逆映射为
		\[
		f_{A^{-1}}:x\longmapsto A^{-1}x.
		\]
		若 $A$ 可逆，对方程 $Ax=\beta$，其有唯一一解
		\[
		x=A^{-1}\beta.
		\]
		
		\item 若 $A$ 有零行，则 $A$ 必不可逆。因为 $\forall B\in M_n(\F)$，$AB$ 有一行全为 $0$，不可能等于 $E_n$。
		
		\item $aE_n$ 可逆当且仅当 $a\neq0$，且
		\[
		(aE_n)^{-1}=a^{-1}E_n.
		\]
		
		\item 根据 \ref{3.3.11}，
		\[
		P(i,j)^{-1}=P(i,j),
		\]
		\[
		P(a(i)+(j))^{-1}=P(-a(i)+(j)),
		\]
		\[
		P(c(i))^{-1}=P(c^{-1}(i)).
		\]
		
		\item 若 $A\in M_n(\F)$ 可逆，则 $r(A)=n$。由
		\[
		n=r(E_n)=r(AA^{-1})\leq r(A)\leq n
		\]
		即可。
	\end{enumerate}
\end{example}


\begin{proposition}\label{3.4.3}
	设 $A$ 可逆,则$A$ 可逆当且仅当 $A$ 的极小多项式$m(x)$的常数项不等于 $0$，且 $A^{-1}$ 是 $A$ 的多项式。
\end{proposition}

\begin{proof}
	由 $m(A)=0$，有
	\[
	0=m(A)=A^k+a_{k-1}A^{k-1}+\cdots+a_1A+a_0E_n.
	\]
	若 $a_0=0$，则
	\[
	A\left(A^{k-1}+a_{k-1}A^{k-2}+\cdots+a_1E_n\right)=0.
	\]
	此时由消去律，得
	\[
	A^{k-1}+a_{k-1}A^{k-2}+\cdots+a_1E_n=0,
	\]
	矛盾，故 $a_0\neq0$。
	于是
	\[
	A\left(A^{k-1}+a_{k-1}A^{k-2}+\cdots+a_1E_n\right)=-a_0E_n,
	\]
	即
	\[
	A\left(-\frac1{a_0}\left(A^{k-1}+a_{k-1}A^{k-2}+\cdots+a_1E_n\right)\right)=E_n.
	\]
	故
	\[
	A^{-1}=-\frac1{a_0}\left(A^{k-1}+a_{k-1}A^{k-2}+\cdots+a_1E_n\right).
	\]
\end{proof}

\begin{proposition}\label{3.4.4}
	设 $A,B\in M_n(\F)$ 且 $A,B$ 可逆，则
	\begin{enumerate}
		\item $A+B$ 不一定可逆；
		\item $A'$ 可逆，且
		\[
		(A')^{-1}=(A^{-1})';
		\]
		\item $AB,BA$ 可逆，且
		\[
		(AB)^{-1}=B^{-1}A^{-1},\qquad (BA)^{-1}=A^{-1}B^{-1};
		\]
		\item $(A^{-1})^{-1}=A$。
	\end{enumerate}
\end{proposition}

\begin{proof}
	$1.$ 令
	\[
	A=\begin{pmatrix}1&1\\0&1\end{pmatrix},
	\qquad
	B=\begin{pmatrix}-1&-1\\0&1\end{pmatrix}.
	\]
	则
	\[
	A+B=\begin{pmatrix}0&0\\0&2\end{pmatrix}
	\]
	不可逆。
	
	$2.$ 由
	\[
	A^{-1}A=AA^{-1}=E_n
	\]
	可得
	\[
	(A^{-1}A)'=(AA^{-1})'=E_n,
	\]
	即
	\[
	A'(A^{-1})'=(A^{-1})'A'=E_n.
	\]
	故
	\[
	(A')^{-1}=(A^{-1})'.
	\]
	
	$3.$ 有
	\[
	(AB)(B^{-1}A^{-1})=A(BB^{-1})A^{-1}=E_n,
	\]
	且
	\[
	(B^{-1}A^{-1})(AB)=B^{-1}(A^{-1}A)B=E_n.
	\]
	所以 $(AB)^{-1}=B^{-1}A^{-1}$。同理
	\[
	(BA)^{-1}=A^{-1}B^{-1}.
	\]
	
	$4.$ 显然。
\end{proof}

\begin{definition}\label{3.4.5}\phantomsection\label{p1:concept:c3-l1171}
	把所有可逆阵构成的集合记为 $\operatorname{GL}_n(\F)$，即
	\[
	\operatorname{GL}_n(\F):=\{A\in M_n(\F):A\text{ 可逆}\}.
	\]
\end{definition}

\begin{proposition}\label{3.4.6}
	对 $A\in M_n(\F)$，$r(A)=n$ 当且仅当 $A$ 可逆。
\end{proposition}

\begin{proof}
	必要性：若 $A$ 可逆，则
	\[
	n=r(E_n)=r(AA^{-1})\leq r(A)\leq n.
	\]
	故 $r(A)=n$。
	
	充分性：由 $n=r(A)$，得 $E_n$ 与 $A$ 相抵。存在
	\[
	P=P_1\cdots P_s,
	\qquad
	Q=Q_1\cdots Q_t,
	\]
	其中 $P_1,\cdots,P_s$ 与 $Q_1,\cdots,Q_t$ 是初等阵，使得
	\[
	PA=AQ=E_n.
	\]
	那么
	\[
	P=PE_n=PAQ=E_nQ=Q.
	\]
	故存在 $B$，使得
	\[
	BA=AB=E_n.
	\]
\end{proof}

\begin{corollary}\label{3.4.7}
	$A\in M_n(\F)$，则 $A$ 可逆当且仅当存在 $B\in M_n(\F)$，使得
	\[
	AB=E_n\quad\text{或}\quad BA=E_n.
	\]
\end{corollary}

\begin{proposition}\label{3.4.8}
	$f,g\in\F[x]$ 且 $(f,g)=1$。设 $f$ 是 $A$ 的零化多项式，则 $g(A)$ 可逆，并且 $g(A)^{-1}$ 是 $A$ 的多项式。
\end{proposition}

\begin{proof}
	由 $(f,g)=1$，存在 $u,v\in\F[x]$，使得
	\[
	uf+vg=1.
	\]
	作用于 $A$，有
	\[
	u(A)f(A)+v(A)g(A)=E_n.
	\]
	又 $f(A)=0$，故
	\[
	g(A)^{-1}=v(A).
	\]
	从而 $g(A)$ 可逆且其逆是 $A$ 的多项式。
\end{proof}

\begin{proposition}\label{3.4.9}
	$s$ 个 $n$ 阶方阵 $A_i(1\leq i\leq s)$ 的乘积可逆，当且仅当每一个 $A_i$ 可逆。
\end{proposition}

\begin{proof}
	必要性：不然，至少存在一个 $A_k$ 不可逆，由 \ref{3.4.3}(3) 矛盾。
	
	充分性：对 \ref{3.4.3}(3) 使用归纳法即可。
\end{proof}

利用 \ref{3.4.2} 中的 4，我们可以得到以下推论。

\begin{corollary}\label{3.4.10}
	$A$ 可逆当且仅当 $A$ 可以写成 $\F$ 上若干个初等阵的乘积。
\end{corollary}

这样，我们就可以得到矩阵求逆的 Gauss--Jordan 算法。

\begin{algo}
	对 $A\in\operatorname{GL}_n(\F)$，存在初等阵 $P_1,\cdots,P_s$，使得
	\[
	P_1\cdots P_sA=E_n.
	\]
	那么
	\[
	A^{-1}=P_1\cdots P_s.
	\]
	也就是说，我们对 $A$ 与 $E_n$ 做一样的初等行变换，使 $A\to E_n$ 时 $E_n\to A^{-1}$，也即
	\[
	(A\ E_n)\longrightarrow (E_n\ A^{-1}).
	\]
	用初等列变换，有
	\[
	\left(\begin{array}{c}
		A\\ 
		E_n
	\end{array}\right)
	\longrightarrow
	\left(\begin{array}{c}
		E_n\\ 
		A^{-1}
	\end{array}\right).
	\]
\end{algo}


\begin{corollary}\label{3.4.12}
	左乘或右乘可逆阵不改变矩阵的秩。
	若 $A\in M_{m\times n}(\F)$，$P\in\operatorname{GL}_m(\F)$，$Q\in\operatorname{GL}_n(\F)$，则
	\[
	r(PA)=r(A)=r(AQ).
	\]
\end{corollary}

\begin{proof}
	有
	\[
	r(PA)\leq r(A)=r(P^{-1}PA)\leq r(PA),
	\]
	从而 $r(PA)=r(A)$。同理
	\[
	r(AQ)\leq r(A)=r(AQQ^{-1})\leq r(AQ),
	\]
	从而 $r(AQ)=r(A)$。
\end{proof}

\begin{corollary}\label{3.4.13}
	$M,N\in M_{m\times n}(\F)$，则 $M$ 与 $N$ 相抵当且仅当存在
	$P\in\operatorname{GL}_m(\F)$，$Q\in\operatorname{GL}_n(\F)$，使得
	\[
	M=PNQ.
	\]
	特别地，$M\in M_{m\times n}(\F)$，$r(M)=r$，当且仅当存在
	$P\in\operatorname{GL}_m(\F)$，$Q\in\operatorname{GL}_n(\F)$，使得
	\[
	PMQ=\begin{pmatrix}
		E_r&0\\
		0&0
	\end{pmatrix}.
	\]
\end{corollary}



\section{分块矩阵}

\begin{definition}\label{3.5.1}\phantomsection\label{p1:concept:c3-l1322}
	称 $s\times t$ 型表格
	\[
	A=
	\begin{pmatrix}
		A_{11}&\cdots&A_{1t}\\
		\vdots&\ddots&\vdots\\
		A_{s1}&\cdots&A_{st}
	\end{pmatrix}
	\]
	是一个 $s\times t$ 分块矩阵。对每个 $i$，$A_{i1},\cdots,A_{it}$ 行数相同；对每个 $j$，$A_{1j},\cdots,A_{sj}$ 列数相同。$A_{ij}$ 称为 $A$ 的 $(i,j)$-块。
	
	对 $A$ 与 $B$，$A=B$ 当且仅当
	\[
	A_{ij}=B_{ij}.
	\]
	定义准上三角阵为
	\[
	\begin{pmatrix}
		A_{11}&A_{12}&\cdots&A_{1s}\\
		0&A_{22}&\cdots&A_{2s}\\
		\vdots&\vdots&\ddots&\vdots\\
		0&0&\cdots&A_{ss}
	\end{pmatrix},
	\]
	准对角阵为
	\[
	\begin{pmatrix}
		A_{11}&&&\\
		&A_{22}&&\\
		&&\ddots&\\
		&&&A_{ss}
	\end{pmatrix}.
	\]
\end{definition}

\begin{proposition}\label{3.5.2}
	\begin{enumerate}
		\item
		\[
		r\begin{pmatrix}
			A&0\\
			0&B
		\end{pmatrix}=r(A)+r(B).
		\]
		\item
		\[
		r\begin{pmatrix}
			A&C\\
			0&B
		\end{pmatrix}\geq r(A)+r(B).
		\]
	\end{enumerate}
\end{proposition}

\begin{proof}
	对 $1$，考虑 $A$ 的列向量组 $\alpha_1,\cdots,\alpha_n$ 的伸长组
	\[
	\begin{pmatrix}\alpha_1\\0\end{pmatrix},\cdots,
	\begin{pmatrix}\alpha_n\\0\end{pmatrix},
	\]
	以及 $B$ 的列向量组 $\beta_1,\cdots,\beta_m$ 的伸长组
	\[
	\begin{pmatrix}0\\\beta_1\end{pmatrix},\cdots,
	\begin{pmatrix}0\\\beta_m\end{pmatrix}.
	\]
	若 $\alpha_{i_1},\cdots,\alpha_{i_r}$ 是 $\alpha_1,\cdots,\alpha_n$ 的极大无关组，
	$\beta_{j_1},\cdots,\beta_{j_q}$ 是 $\beta_1,\cdots,\beta_m$ 的极大无关组，则
	\[
	\begin{pmatrix}\alpha_{i_1}\\0\end{pmatrix},\cdots,
	\begin{pmatrix}\alpha_{i_r}\\0\end{pmatrix},
	\begin{pmatrix}0\\\beta_{j_1}\end{pmatrix},\cdots,
	\begin{pmatrix}0\\\beta_{j_q}\end{pmatrix}
	\]
	无关，故
	\[
	r\begin{pmatrix}
		A&0\\0&B
	\end{pmatrix}=r(A)+r(B).
	\]
	
	对 $2$，考虑
	\[
	\begin{pmatrix}
		A&C\\0&B
	\end{pmatrix}.
	\]
	取 $A$ 的列向量组中一个极大无关组 $\alpha_{i_1},\cdots,\alpha_{i_r}$，
	取 $B$ 的列向量组中一个极大无关组 $\beta_{j_1},\cdots,\beta_{j_q}$。若 $C$ 的相应列记为
	$\gamma_{j_1},\cdots,\gamma_{j_q}$，
	\[
	\begin{pmatrix}\alpha_{i_1}\\0\end{pmatrix},\cdots,
	\begin{pmatrix}\alpha_{i_r}\\0\end{pmatrix},
	\begin{pmatrix}\gamma_{j_1}\\\beta_{j_1}\end{pmatrix},\cdots,
	\begin{pmatrix}\gamma_{j_q}\\\beta_{j_q}\end{pmatrix}
	\]
	可能存在$\gamma_i$使$ \alpha_{i_1},\cdots,\alpha_{i_r},\gamma_i$无关，故
	\[
	r\begin{pmatrix}
		A&C\\0&B
	\end{pmatrix}\geq r(A)+r(B).
	\]
\end{proof}

一个常用的分块方式是将 $(1,1)$-元或 $(m,n)$-元单分一行，也即
\[
A=\begin{pmatrix}
	a_{11}&\alpha'\\
	\beta&A_1
\end{pmatrix}
=\begin{pmatrix}
	A'_1&\alpha\\
	\beta'&a_{mn}
\end{pmatrix}.
\]

分块矩阵的运算定义如下。

\begin{definition}\label{3.5.3}\phantomsection\label{p1:concept:c3-l1440}
	\begin{enumerate}
		\item 转置：
		\[
		\begin{pmatrix}
			A_{11}&\cdots&A_{1t}\\
			\vdots&\ddots&\vdots\\
			A_{s1}&\cdots&A_{st}
		\end{pmatrix}'
		=
		\begin{pmatrix}
			A_{11}'&\cdots&A_{s1}'\\
			\vdots&\ddots&\vdots\\
			A_{1t}'&\cdots&A_{st}'
		\end{pmatrix}.
		\]
		
		\item 数乘：
		\[
		k\begin{pmatrix}
			A_{11}&\cdots&A_{1t}\\
			\vdots&\ddots&\vdots\\
			A_{s1}&\cdots&A_{st}
		\end{pmatrix}
		=
		\begin{pmatrix}
			kA_{11}&\cdots&kA_{1t}\\
			\vdots&\ddots&\vdots\\
			kA_{s1}&\cdots&kA_{st}
		\end{pmatrix}.
		\]
		
		\item 加法：对 $A=(A_{ij})$，$B=(B_{ij})$，分块的方式保证 $A_{ij}+B_{ij}$ 对任意 $i,j$ 有定义，则
		\[
		\begin{pmatrix}
			A_{11}&\cdots&A_{1t}\\
			\vdots&\ddots&\vdots\\
			A_{s1}&\cdots&A_{st}
		\end{pmatrix}
		+
		\begin{pmatrix}
			B_{11}&\cdots&B_{1t}\\
			\vdots&\ddots&\vdots\\
			B_{s1}&\cdots&B_{st}
		\end{pmatrix}
		=
		\begin{pmatrix}
			A_{11}+B_{11}&\cdots&A_{1t}+B_{1t}\\
			\vdots&\ddots&\vdots\\
			A_{s1}+B_{s1}&\cdots&A_{st}+B_{st}
		\end{pmatrix}.
		\]
		
		\item 乘法：对 $C=AB$，$A=(A_{ij})_{s,t}$，$B=(B_{ij})_{t,r}$，$C=(C_{ij})_{s,r}$，
		\[
		\begin{pmatrix}
			A_{11}&\cdots&A_{1t}\\
			\vdots&\ddots&\vdots\\
			A_{s1}&\cdots&A_{st}
		\end{pmatrix}
		\begin{pmatrix}
			B_{11}&\cdots&B_{1r}\\
			\vdots&\ddots&\vdots\\
			B_{t1}&\cdots&B_{tr}
		\end{pmatrix}
		=
		\begin{pmatrix}
			C_{11}&\cdots&C_{1r}\\
			\vdots&\ddots&\vdots\\
			C_{s1}&\cdots&C_{sr}
		\end{pmatrix},
		\]
		其中
		\[
		C_{pq}=A_{p1}B_{1q}+\cdots+A_{pt}B_{tq}.
		\]
		$A,B$ 的分块方式保证所有的矩阵乘法与加法有定义。
	\end{enumerate}
\end{definition}

那么，我们也可以定义出分块矩阵上的初等变换。

根据上面的定义，不难得出下面的推论。

\begin{corollary}\label{3.5.4}
	\[
	\begin{pmatrix}
		A_{11}&A_{12}&\cdots&A_{1s}\\
		0&A_{22}&\cdots&A_{2s}\\
		\vdots&\vdots&\ddots&\vdots\\
		0&0&\cdots&A_{ss}
	\end{pmatrix}^{k}
	=
	\begin{pmatrix}
		A_{11}^{k}&*&\cdots&*\\
		0&A_{22}^{k}&\cdots&*\\
		\vdots&\vdots&\ddots&\vdots\\
		0&0&\cdots&A_{ss}^{k}
	\end{pmatrix}.
	\]
	可以看到，准上三角阵的任意多项式仍是准上三角阵。
\end{corollary}
下面可以看到一些使用分块矩阵的应用
\begin{proposition}\label{3.5.5}
	对 $A\in M_{m\times n}(\F)$，$B\in M_{n\times s}(\F)$ 且 $AB=0$，则
	\[
	r(A)+r(B)\leq n.
	\]
\end{proposition}

\begin{proof}
	记
	\[
	B=(\eta_1\ \cdots\ \eta_s).
	\]
	则
	\[
	AB=(A\eta_1\ \cdots\ A\eta_s)=0,
	\]
	也即
	\[
	\left\{
	\begin{aligned}
		A\eta_1&=0,\\
		&\vdots\\
		A\eta_s&=0.
	\end{aligned}
	\right.
	\]
	那么 $\eta_i(1\leq i\leq s)$ 是 $Ax=0$ 的解。任取 $Ax=0$ 的基础解系，取其中 $n-r(A)$ 个无关的解向量
	\[
	\xi_1,\cdots,\xi_{n-r(A)}.
	\]
	$\eta_1,\cdots,\eta_s$ 可由 $\xi_1,\cdots,\xi_{n-r(A)}$ 表出，由基本引理，
	\[
	r(\eta_1,\cdots,\eta_s)=r(B)\leq n-r(A),
	\]
	故
	\[
	r(A)+r(B)\leq n.
	\]
\end{proof}

\begin{proposition}\label{3.5.6}
	$f,g\in\F[x]$ 且 $(f,g)=1$，$A\in M_n(\F)$，则
	\[
	f(A)g(A)=0
	\]
	当且仅当
	\[
	r(f(A))+r(g(A))=n.
	\]
\end{proposition}

\begin{proof}
	由 $(f,g)=1$，存在 $u,v\in\F[x]$，使得
	\[
	fu+gv=1.
	\]
	代入 $A$，有
	\[
	f(A)u(A)+g(A)v(A)=E_n.
	\]
	考虑以下初等变换
	\[
	\begin{pmatrix}
		f(A)&0\\
		0&g(A)
	\end{pmatrix}
	\longrightarrow
	\begin{pmatrix}
		f(A)&f(A)u(A)\\
		0&g(A)
	\end{pmatrix}
	\longrightarrow
	\begin{pmatrix}
		f(A)&E_n\\
		0&g(A)
	\end{pmatrix}
	\]
	\[
	\longrightarrow
	\begin{pmatrix}
		f(A)&E_n\\
		-f(A)g(A)&0
	\end{pmatrix}
	\longrightarrow
	\begin{pmatrix}
		0&E_n\\
		-f(A)g(A)&0
	\end{pmatrix}
	\longrightarrow
	\begin{pmatrix}
		f(A)g(A)&0\\
		0&E_n
	\end{pmatrix}.
	\]
	故有
	\[
	r(f(A))+r(g(A))=r(f(A)g(A))+n,
	\]
	从而得证。
\end{proof}

在上一章中，我们给出 $Ax=\beta$ 有解的必要条件是
\[
r(A\ \beta)=r(A).
\]
对矩阵方程 $AX=B$，其有解的必要条件是
\[
r(A\ B)=r(A).
\]

\begin{proposition}\label{3.5.7}
	$A\in M_{m\times n}(\F)$，$B\in M_{m\times s}(\F)$，则存在 $X\in M_{n\times s}(\F)$ 是 $AX=B$ 的解的必要条件是
	\[
	r(A\ B)=r(A).
	\]
\end{proposition}

\begin{proof}
	记
	\[
	A=(\alpha_1\ \cdots\ \alpha_n),
	\qquad
	B=(\eta_1\ \cdots\ \eta_s).
	\]
	那么
	\[
	(\alpha_1\ \cdots\ \alpha_n)
	\begin{pmatrix}
		x_{11}&\cdots&x_{1s}\\
		\vdots&\ddots&\vdots\\
		x_{n1}&\cdots&x_{ns}
	\end{pmatrix}
	=(\eta_1\ \cdots\ \eta_s).
	\]
	于是
	\[
	\eta_i=x_{1i}\alpha_1+\cdots+x_{ni}\alpha_n,
	\]
	也即 $\eta_1,\cdots,\eta_s$ 可由 $\alpha_1,\cdots,\alpha_n$ 表出，故
	\[
	r(A\ B)=r(A).
	\]
\end{proof}

\begin{lemma}\label{3.5.8}
	$\forall A\in M_{m\times n}(\mathbb R)$，则有
	\[
	r(A'A)=r(A).
	\]
\end{lemma}

\begin{proof}
	考虑
	\[
	\text{(I)}\quad Ax=0,
	\]
	其基础解系中含 $n-r(A)$ 个向量；以及
	\[
	\text{(II)}\quad A'Ax=0,
	\]
	其基础解系中含 $n-r(A'A)$ 个向量。
	显然，(I) 的解都是 (II) 的解。下证 (II) 的解都是 (I) 的解。
	
	对 $\xi\in\mathbb R^n$ 且 $A'A\xi=0$，证明 $A\xi=0$。
	由 $A'A\xi=0$，左乘 $\xi'$，得
	\[
	\xi'A'A\xi=0\quad\Longleftrightarrow\quad (A\xi)'A\xi=0.
	\]
	设
	\[
	A\xi=(a_1\ \cdots\ a_m)'\in\mathbb R^m,
	\]
	则
	\[
	a_1^2+\cdots+a_m^2=0,
	\qquad a_1,\cdots,a_m\in\mathbb R.
	\]
	故 $a_1=\cdots=a_m=0$，也即 $A\xi=0$。
	综上，(I)、(II) 的解空间相同，故
	\[
	n-r(A)=n-r(A'A),
	\]
	即
	\[
	r(A)=r(A'A).
	\]
\end{proof}

\begin{proposition}\label{3.5.9}
	$A\in M_{m\times n}(\mathbb R)$，$\forall\beta\in\mathbb R^m$，方程
	\[
	A'Ax=A'\beta
	\]
	总是有解。
\end{proposition}

\begin{proof}
	只需证
	\[
	r(A'A)=r(A'A\ ~~ A'\beta),
	\]
	而
	\[
	(A'A\ ~~ A'\beta)=A'(A\ \beta),
	\]
	故
	\[
	r(A'A)\leq r(A'A\ ~~ A'\beta)
	= r\bigl(A'(A\ \beta)\bigr)
	\leq r(A')=r(A)=r(A'A).
	\]
	得证。
\end{proof}

对分块矩阵
\[
\begin{pmatrix}
	A&B\\
	C&D
\end{pmatrix},
\]
若 $A\in\operatorname{GL}_m(\F)$，考虑其上的 Gauss 消元法：
\[
\begin{pmatrix}
	E_m&0\\
	-CA^{-1}&E_n
\end{pmatrix}
\begin{pmatrix}
	A&B\\
	C&D
\end{pmatrix}
=
\begin{pmatrix}
	A&B\\
	0&D-CA^{-1}B
\end{pmatrix},
\]
\[
\begin{pmatrix}
	A&B\\
	C&D
\end{pmatrix}
\begin{pmatrix}
	E_m&-A^{-1}B\\
	0&E_n
\end{pmatrix}
=
\begin{pmatrix}
	A&0\\
	C&D-CA^{-1}B
\end{pmatrix}.
\]

使用这个算法，也可以推出准上三角阵的逆。若
\[
M=\begin{pmatrix}
	A&B\\
	0&C
\end{pmatrix},
\]
且 $A,C$ 可逆，则 $M$ 可逆且
\[
M^{-1}=\begin{pmatrix}
	A^{-1}&-A^{-1}BC^{-1}\\
	0&C^{-1}
\end{pmatrix}.
\]
\section{行列式}

\begin{definition}\label{3.6.1}\phantomsection\label{p1:concept:c3-l1812}
	把任意 $n$ 个给定的互异的对象按任意顺序排成一列，就得到这 $n$ 个对象的一个排列，称为一个 $n$ 元排列。
	
	把这 $n$ 个不同的对象编号为 $1,2,\cdots,n$，因此一个 $n$ 元排列可以写成
	\[
	i_1i_2\cdots i_n,
	\]
	其中 $i_k$ 表示第 $k$ 个位置上的数。
	排列 $12\cdots n$ 称为自然排列。
	
	$S_n$ 表示所有 $n$ 元排列构成的集合，$S_n$ 称为 $n$ 元置换群，且
	\[
	|S_n|=n!.
	\]
	$S_n$ 是一个有限群。
\end{definition}
实际上，$i_1\cdots i_n$ 可以看作一个双射，也即置换：
\[
i_1\cdots i_n
\Longleftrightarrow
\begin{pmatrix}
	1&2&\cdots&n\\
	i_1&i_2&\cdots&i_n
\end{pmatrix}.
\]
该置换表示把
\[
1\mapsto i_1,
\quad
2\mapsto i_2,
\quad\cdots\quad,
n\mapsto i_n.
\]
\begin{definition}\label{3.6.2}\phantomsection\label{p1:concept:c3-l1845}
	对 $i_1\cdots i_n\in S_n$，若 $j<k$ 但 $i_j>i_k$，称 $(i_j,i_k)$ 是一个逆对。
	$i_1\cdots i_n$ 所有逆对的个数称为逆序数，记为
	\[
	\tau(i_1\cdots i_n).
	\]
	称
	\[
	(-1)^{\tau(i_1\cdots i_n)}
	\]
	为排列 $i_1\cdots i_n$ 的符号。
	
	若 $\tau(i_1\cdots i_n)$ 为奇数，则称 $i_1\cdots i_n$ 为奇排列；若 $\tau(i_1\cdots i_n)$ 为偶数，则称 $i_1\cdots i_n$ 为偶排列。
	
\end{definition}
例如
\[
\tau(12\cdots n)=0,
\qquad
\tau\bigl(n(n-1)\cdots 21\bigr)=C_n^2,
\]
并且
\[
\tau(31542)=2+0+2+1=5.
\]
\begin{definition}\label{3.6.3}\phantomsection\label{p1:concept:c3-l1870}
	对 $1\leq j,k\leq n$，定义一个从 $S_n$ 到 $S_n$ 的映射，称为对换：
	\[
	(j\ k):
	(i_1\cdots i_j\cdots i_k\cdots i_n)
	\longmapsto
	(i_1\cdots i_k\cdots i_j\cdots i_n).
	\]
\end{definition}

\begin{lemma}\label{3.6.4}
	对换改变一个排列的奇偶性。
\end{lemma}

\begin{proof}
	设 $k=j+s$。先考虑相邻对换 $(j\ j+1)$，显然相邻两项交换会改变逆对个数的奇偶性。
	再看 $(j\ j+2)$，它可看成三次相邻对换的复合：
	\[
	(j\ j+2)=(j\ j+1)(j+1\ j+2)(j\ j+1).
	\]
	一般地，由归纳可知
	\[
	(j\ j+s)
	\]
	是 $2s-1$ 次相邻对换的复合。由于 $2s-1$ 为奇数，故任意对换都会改变排列的奇偶性。
\end{proof}

\begin{lemma}\label{3.6.5}
	$\forall i_1\cdots i_n\in S_n$，可以经过若干次对换得到 $12\cdots n$。
\end{lemma}

\begin{proof}
	使用归纳法。$n=2$ 时显然。设 $n=m$ 时结论成立。
	若在排列 $i_1\cdots i_m$ 中有 $i_k=m$，则作对换 $(k\ m)$，把 $m$ 换到最后一个位置。
	然后对剩下的 $m-1$ 个数由归纳假设可经若干次对换得到自然排列。
	故任意排列均可经过若干次对换得到 $12\cdots n$。
\end{proof}

\begin{definition}\label{3.6.6}\phantomsection\label{p1:concept:c3-l1908}
	给定域 $\F$。对任意 $A\in M_n(\F)$，$A=(a_{ij})$，任选 $j_1\cdots j_n\in S_n$，考虑
	\[
	a_{1j_1}a_{2j_2}\cdots a_{nj_n}\in\F,
	\]
	就有
	\[
	(-1)^{\tau(j_1\cdots j_n)}a_{1j_1}a_{2j_2}\cdots a_{nj_n}
	\]
	一共 $n!$ 个这样的项。定义 $A$ 的行列式 $\det A$ 为
	\[
	\det A=
	\sum_{j_1\cdots j_n\in S_n}
	(-1)^{\tau(j_1\cdots j_n)}a_{1j_1}a_{2j_2}\cdots a_{nj_n}.
	\]
	也记为
	\[
	|A|,
	\qquad
	\det(a_{ij}),
	\qquad
	\left|
	\begin{matrix}
		a_{11}&\cdots&a_{1n}\\
		\vdots&\ddots&\vdots\\
		a_{n1}&\cdots&a_{nn}
	\end{matrix}
	\right|,
	\]
	称为 $n$ 阶行列式。
\end{definition}

下面给出几6个关于行列式的例子。

\begin{example}
	对 $f_1,\cdots,f_n\in D^{n-1}(a,b)$，考虑 $f_1,\cdots,f_n$ 的线性相关性：
	\[
	\left\{
	\begin{aligned}
		k_1f_1+\cdots+k_nf_n&=0,\\
		k_1f_1'+\cdots+k_nf_n'&=0,\\
		&\vdots\\
		k_1f_1^{(n-1)}+\cdots+k_nf_n^{(n-1)}&=0.
	\end{aligned}
	\right.
	\]
	设
	\[
	W(f_1,\cdots,f_n)=
	\left|
	\begin{matrix}
		f_1&f_2&\cdots&f_n\\
		f_1'&f_2'&\cdots&f_n'\\
		\vdots&\vdots&\ddots&\vdots\\
		f_1^{(n-1)}&f_2^{(n-1)}&\cdots&f_n^{(n-1)}
	\end{matrix}
	\right|.
	\]
	后面我们将看到 $W(f_1,\cdots,f_n)$ 的值决定了 $f_1,\cdots,f_n$ 是否相关。
\end{example}


\begin{example}
	设 $A=(a_{ij})\in M_n(\F)$，$\lambda$ 是不定元，令
	\[
	\begin{aligned}
		f(\lambda)
		&=|\lambda E_n-A|\\
		&=
		\left|
		\begin{matrix}
			\lambda-a_{11}&-a_{12}&\cdots&-a_{1n}\\
			-a_{21}&\lambda-a_{22}&\cdots&-a_{2n}\\
			\vdots&\vdots&\ddots&\vdots\\
			-a_{n1}&-a_{n2}&\cdots&\lambda-a_{nn}
		\end{matrix}
		\right|.
	\end{aligned}
	\]
	\phantomsection\label{p1:concept:c3-l1987}$f(\lambda)$ 是 $\F$ 上的首一的 $n$ 次多项式，$f(\lambda)$ 称为 $A$ 的特征多项式。
\end{example}

\begin{example}
	对 $S_1=\{1\}$，有
	\[
	\det(a)=a.
	\]
	对 $S_2=\{12,21\}$，有
	\[
	\begin{aligned}
		\det
		\begin{pmatrix}
			a_{11}&a_{12}\\
			a_{21}&a_{22}
		\end{pmatrix}
		&=(-1)^{\tau(12)}a_{11}a_{22}
		+(-1)^{\tau(21)}a_{12}a_{21}\\
		&=a_{11}a_{22}-a_{12}a_{21}.
	\end{aligned}
	\]
	$n=3$ 时，
	\[
	\begin{aligned}
		\det
		\begin{pmatrix}
			a_{11}&a_{12}&a_{13}\\
			a_{21}&a_{22}&a_{23}\\
			a_{31}&a_{32}&a_{33}
		\end{pmatrix}
		=&(-1)^{\tau(123)}a_{11}a_{22}a_{33}
		+(-1)^{\tau(132)}a_{11}a_{23}a_{32}\\
		&+(-1)^{\tau(213)}a_{12}a_{21}a_{33}
		+(-1)^{\tau(231)}a_{12}a_{23}a_{31}\\
		&+(-1)^{\tau(312)}a_{13}a_{21}a_{32}
		+(-1)^{\tau(321)}a_{13}a_{22}a_{31}.
	\end{aligned}
	\]
\end{example}


若 $A$ 有零行或零列，则 $|A|=0$。反之不一定，例如
\[
A=\begin{pmatrix}1&1\\1&1\end{pmatrix},
\qquad
|A|=0.
\]

对上三角矩阵，
\[
\left|
\begin{matrix}
	a_{11}&a_{12}&\cdots&a_{1n}\\
	0&a_{22}&\cdots&a_{2n}\\
	\vdots&\vdots&\ddots&\vdots\\
	0&0&\cdots&a_{nn}
\end{matrix}
\right|
=a_{11}a_{22}\cdots a_{nn}.
\]

\section{行列式的性质}

\begin{lemma}\label{3.7.1}
	设 $A=(a_{ij})\in M_n(\F)$，$i_1\cdots i_n\in S_n$，则
	\[
	(-1)^{\tau(i_1\cdots i_n)+\tau(j_1\cdots j_n)}
	a_{i_1j_1}\cdots a_{i_nj_n}
	\]
	是 $\det A$ 的一个和项。
\end{lemma}

\begin{proof}
	对 $i_1\cdots i_n$，经过 $s$ 次对换后变为 $1\cdots n$，则
	\[
	(-1)^s=(-1)^{\tau(i_1\cdots i_n)}.
	\]
	设 $j_1\cdots j_n$ 经过相应的 $s$ 次对换后变为 $k_1\cdots k_n$，则
	\[
	(-1)^s(-1)^{\tau(j_1\cdots j_n)}=(-1)^{\tau(k_1\cdots k_n)}.
	\]
	由乘法交换律得
	\[
	a_{i_1j_1}\cdots a_{i_nj_n}=a_{1k_1}\cdots a_{nk_n},
	\]
	且
	\[
	(-1)^{\tau(i_1\cdots i_n)+\tau(j_1\cdots j_n)}
	=(-1)^s(-1)^{\tau(j_1\cdots j_n)}
	=(-1)^{\tau(k_1\cdots k_n)}.
	\]
	故该项是 $\det A$ 的一个和项。
\end{proof}

\begin{proposition}\label{3.7.2}
	设 $A=(a_{ij})\in M_n(\F)$。
	\begin{enumerate}
		\item 对任意 $i_1\cdots i_n\in S_n$，
		\[
		|A|=(-1)^{\tau(i_1\cdots i_n)}
		\sum_{j_1\cdots j_n\in S_n}
		(-1)^{\tau(j_1\cdots j_n)}
		a_{i_1j_1}\cdots a_{i_nj_n}.
		\]
		\item 对任意 $j_1\cdots j_n\in S_n$，
		\[
		|A|=(-1)^{\tau(j_1\cdots j_n)}
		\sum_{i_1\cdots i_n\in S_n}
		(-1)^{\tau(i_1\cdots i_n)}
		a_{i_1j_1}\cdots a_{i_nj_n}.
		\]
	\end{enumerate}
\end{proposition}

\begin{proof}
	对 $1$，左右两侧的项数一样，由上引理易得。
	$2$ 同理。
\end{proof}

我们可以得到以下推论。

\begin{corollary}\label{3.7.3}
	\[
	|A|=
	\sum_{i_1\cdots i_n\in S_n}
	(-1)^{\tau(i_1\cdots i_n)}
	a_{i_11}\cdots a_{i_nn}.
	\]
\end{corollary}

\begin{corollary}\label{3.7.4}
	\[
	|A|=|A'|.
	\]
\end{corollary}

下面，考虑初等变换对矩阵行列式的影响，让我们一类一类地考虑对应的矩阵。

\begin{proposition}\label{3.7.5}
	设 $A\in M_n(\F)$，则
	\[
	|P(i,j)A|=-|A|=|AP(i,j)|.
	\]
\end{proposition}

\begin{proof}
	只证 $|P(i,j)A|=-|A|$。设 $P(i,j)A=B$，则
	\[
	\begin{aligned}
		|B|
		&=\sum_{k_1\cdots k_i\cdots k_j\cdots k_n\in S_n}
		(-1)^{\tau(k_1\cdots k_i\cdots k_j\cdots k_n)}
		B(1,k_1)\cdots B(i,k_i)\cdots B(j,k_j)\cdots B(n,k_n)\\
		&=\sum_{k_1\cdots k_i\cdots k_j\cdots k_n\in S_n}
		(-1)^{\tau(k_1\cdots k_i\cdots k_j\cdots k_n)}
		A(1,k_1)\cdots A(j,k_i)\cdots A(i,k_j)\cdots A(n,k_n)\\
		&=\sum_{k_1\cdots k_i\cdots k_j\cdots k_n\in S_n}
		(-1)^{\tau(k_1\cdots k_i\cdots k_j\cdots k_n)}
		\frac{(-1)^{\tau(k_1\cdots k_j\cdots k_i\cdots k_n)}}
		{(-1)^{\tau(k_1\cdots k_j\cdots k_i\cdots k_n)}}\\
		&\qquad\qquad\cdot
		A(1,k_1)\cdots A(i,k_j)\cdots A(j,k_i)\cdots A(n,k_n)\\
		&=-|A|.
	\end{aligned}
	\]
	同理可证 $|AP(i,j)|=-|A|$。
\end{proof}

应用这一命题，不难得出：若有两行或两列相同，则 $|A|=0$。类似于对 $P(i,j)$ 相同的做法，不难有 $P(c(i))$ 的公式。

\begin{proposition}\label{3.7.6}
	\[
	|P(c(i))A|=c|A|=|AP(c(i))|.
	\]
\end{proposition}

\begin{proposition}\label{3.7.7}
	设 $A=(\alpha_1\ \cdots\ \alpha_n)\in M_n(\F)$，且 $\alpha_i=\beta+\gamma$，则
	\[
	\det(\alpha_1\ \cdots\ \alpha_i\ \cdots\ \alpha_n)
	=
	\det(\alpha_1\ \cdots\ \beta\ \cdots\ \alpha_n)
	+
	\det(\alpha_1\ \cdots\ \gamma\ \cdots\ \alpha_n).
	\]
\end{proposition}

由 \ref{3.7.7}，我们可以得到以下推论。

\begin{corollary}\label{3.7.8}
	\[
	|P(c(i)+(j))A|=|AP(c(i)+(j))|=|A|.
	\]
\end{corollary}

可以看到，对初等矩阵 $P$，有
\[
|PA|=|P||A|=|AP|.
\]

\begin{theorem}\label{3.7.9}
	设 $A\in M_n(\F)$，则 $A$ 可逆当且仅当 $|A|\neq0$。
\end{theorem}

\begin{proof}
	必要性：若 $A$ 可逆，则 $r(A)=n$，于是 $A$ 与 $E_n$ 相抵，而 $|E_n|=1$，故 $|A|\neq0$。
	
	充分性：若 $|A|\neq0$。若 $r(A)=r<n$，则 $A$ 与
	\[
	\begin{pmatrix}
		E_r&0\\
		0&0
	\end{pmatrix}
	\]
	相抵，而
	\[
	\left|
	\begin{matrix}
		E_r&0\\
		0&0
	\end{matrix}
	\right|=0.
	\]
	这样推出 $|A|=0$，矛盾，故 $r(A)=n$，$A$ 可逆。
\end{proof}

注意，一般来讲 $|A+B|\neq |A|+|B|$。利用上面的公式，例如
\[
\begin{aligned}
	&\det(a_1\xi_1+a_2\xi_2,
	a_3\xi_3+a_4\xi_4)\\
	&=\det(a_1\xi_1,
	a_3\xi_3+a_4\xi_4)
	+\det(a_2\xi_2,
	a_3\xi_3+a_4\xi_4)\\
	&=\det(a_1\xi_1,a_3\xi_3)
	+\det(a_1\xi_1,a_4\xi_4)
	+\det(a_2\xi_2,a_3\xi_3)
	+\det(a_2\xi_2,a_4\xi_4).
\end{aligned}
\]

\begin{proposition}\label{3.7.10}
	设 $A,B\in M_n(\F)$，则
	\begin{enumerate}
		\item $|kA|=k^n|A|$；
		\item $|AB|=|A||B|$；
		\item $|A^{-1}|=|A|^{-1}$；
		\item $|A^k|=|A|^k$。
	\end{enumerate}
\end{proposition}

\begin{proof}
	$1$ 由 \ref{3.7.6} 可得。
	$3,4$ 由 $2$ 得。这里采用扰动法证明 $2$。
	
	若 $A$ 可逆，则 $A=P_1\cdots P_s$，其中 $P_1,\cdots,P_s$ 是初等阵。那么
	\[
	\begin{aligned}
		|AB|&=|P_1\cdots P_sB|\\
		&=|P_1|\cdots|P_s||B|\\
		&=|P_1\cdots P_s||B|\\
		&=|A||B|.
	\end{aligned}
	\]
	若 $A$ 不可逆，断言存在 $t$，使 $tE_n+A$ 可逆。
	考虑
	\[
	f(t)=|tE_n+A|
	\]
	是一个首一的非零多项式，那么 $f(t)$ 只有有限多个零点，则存在 $t$，使 $|tE_n+A|\neq0$，此时 $tE_n+A$ 可逆。
	取无穷小量 $t$，使 $tE_n+A$ 可逆，从而
	\[
	|(tE_n+A)B|=|tE_n+A||B|.
	\]
	令 $t\to0$，得
	\[
	|AB|=|A||B|.
	\]
\end{proof}

用行列式的语言，重新考虑线性齐次方程组的解，有
\[
Ax=0\quad\text{且}\quad A\in M_n(\F)
\quad\text{有非平凡解}\quad\Longleftrightarrow\quad |A|=0.
\]

\begin{corollary}\label{3.7.11}
	设 $A\in M_{m\times n}(\F)$，$B\in M_{n\times m}(\F)$，且 $m\neq n$，则 $|AB|$ 与 $|BA|$ 中至少有一个是 $0$。
\end{corollary}

\begin{proof}
	若 $m<n$，则 $BA\in M_n(\F)$，并且
	\[
	r(BA)\leq\min\{m,n\}=m<n.
	\]
	故 $BA$ 不满秩，从而 $|BA|=0$。
	$m>n$ 同理。
\end{proof}

定义
\[
SL_n(\F):=\{A\in M_n(\F):\det A=1\},
\]
\phantomsection\label{p1:concept:c3-l2291}称为 $n$ 阶特殊线性群，在抽象代数中进一步介绍。

关于扰动法，这里再指出一例。
\begin{example}
	对任意 $X,Y\in M_n(\F)$，证明
	\[
	|E_n-XY|=|E_n-YX|.
	\]
\end{example} 

\begin{proof}
	若 $X$ 可逆，则
	\[
	E_n-XY=X(X^{-1}-Y),
	\]
	且
	\[
	E_n-YX=(X^{-1}-Y)X.
	\]
	那么
	\[
	|E_n-XY|=|X||X^{-1}-Y|=|X^{-1}-Y||X|=|E_n-YX|.
	\]
	若 $X$ 不可逆，考虑
	\[
	X_t=X+tE_n.
	\]
	总存在 $t$ 使 $X_t$ 可逆，那么
	\[
	|E_n-(X+tE_n)Y|=|E_n-Y(X+tE_n)|.
	\]
	令 $t\to0$，得
	\[
	|E_n-XY|=|E_n-YX|.
	\]
\end{proof}


\section{Laplace 定理}

\begin{definition}\label{3.8.1}\phantomsection\label{p1:concept:c3-l2331}
	设 $A\in M_n(\F)$，$n>1$。对 $A(i,j)$，把 $|A|$ 划去第 $i$ 行、第 $j$ 列留下的 $n-1$ 阶行列式记为 $M_{ij}$，称为 $A(i,j)$ 的余子式；称
	\[
	(-1)^{i+j}M_{ij}
	\]
	为 $A(i,j)$ 的代数余子式，记为 $C_{ij}$。
\end{definition}

\begin{theorem}(Laplace 定理的展开式形式)\label{3.8.2}
	对任意 $A\in M_n(\F)$，有
	\[
	\begin{aligned}
		\text{任取第 }i\text{ 行，}
		\quad |A|&=A(i,1)C_{i1}+\cdots+A(i,n)C_{in},\\
		\text{任取第 }j\text{ 列，}
		\quad |A|&=A(1,j)C_{1j}+\cdots+A(n,j)C_{nj}.
	\end{aligned}
	\]
\end{theorem}

\begin{proof}
	只证第 $i$ 行的情况。由定义，
	\[
	\begin{aligned}
		|A|
		&=\sum_{ j_1\cdots  j_n\in S_n}
		(-1)^{\tau( j_1\cdots j_n)}
		A(1, j_1)\cdots A(i-1, j_{i-1})A(i, j_i)A(i+1, j_{i+1})\cdots A(n, j_n)\\
		&=\sum_{k=1}^{n}A(i,k)
		\sum_{ j_1\cdots j_{i-1} j_{i+1}\cdots j_n\in S_{n,k}}
		(-1)^{\tau( j_1\cdots j_{i-1}k j_{i+1}\cdots j_n)}\\
		&\qquad\qquad\cdot
		A(1, j_1)\cdots A(i-1, j_{i-1})A(i+1, j_{i+1})\cdots A(n, j_n),
	\end{aligned}
	\]
	其中
	\[
	S_{n,k}=\{1,2,\ldots,k-1,k+1,\ldots,n\}
	\]
	是$n-1$ 元排列。
	比较
	\[
	(-1)^{\tau( j_1\cdots j_{i-1}k j_{i+1}\cdots j_n)}
	\quad\text{与}\quad
	(-1)^{\tau( j_1\cdots j_{i-1} j_{i+1}\cdots j_n)}.
	\]
	把 $k$ 对换到首位，需要 $i-1$ 次，故
	\[
	\begin{aligned}
		(-1)^{\tau( j_1\cdots j_{i-1}k j_{i+1}\cdots j_n)}
		&=(-1)^{i-1}(-1)^{\tau(k j_1\cdots j_{i-1} j_{i+1}\cdots j_n)}\\
		&=(-1)^{i-1}(-1)^{k-1}
		(-1)^{\tau( j_1\cdots j_{i-1} j_{i+1}\cdots j_n)}\\
		&=(-1)^{i+k}
		(-1)^{\tau( j_1\cdots j_{i-1} j_{i+1}\cdots j_n)}.
	\end{aligned}
	\]
	这样
	\[
	|A|=\sum_{k=1}^{n}A(i,k)C_{ik}.
	\]
	列的情况同理。
\end{proof}
\begin{example}
	（Vandermonde 行列式）设
	\[
	V(x_1,\cdots,x_n)=
	\begin{vmatrix}
		1&1&\cdots&1\\
		x_1&x_2&\cdots&x_n\\
		x_1^2&x_2^2&\cdots&x_n^2\\
		\vdots&\vdots&\ddots&\vdots\\
		x_1^{n-1}&x_2^{n-1}&\cdots&x_n^{n-1}
	\end{vmatrix}.
	\]当 $n>2$ 时，
	\[
	V(x_1,\cdots,x_n)\neq0
	\quad\Longleftrightarrow\quad
	x_1,\cdots,x_n\text{ 互不相同}.
	\]
	事实上，
	\[
	\begin{aligned}
		V(x_1,\cdots,x_n)
		&=
		\begin{vmatrix}
			1&1&\cdots&1\\
			x_1&x_2&\cdots&x_n\\
			\vdots&\vdots&\ddots&\vdots\\
			x_1^{n-1}&x_2^{n-1}&\cdots&x_n^{n-1}
		\end{vmatrix}\\
		&=
		\begin{vmatrix}
			1&1&\cdots&1\\
			0&x_2-x_1&\cdots&x_n-x_1\\
			\vdots&\vdots&\ddots&\vdots\\
			0&x_2^{n-1}-x_1x_2^{n-2}&\cdots&x_n^{n-1}-x_1x_n^{n-2}
		\end{vmatrix}\\
		&=
		\begin{vmatrix}
			x_2-x_1&x_3-x_1&\cdots&x_n-x_1\\
			x_2^2-x_1x_2&x_3^2-x_1x_3&\cdots&x_n^2-x_1x_n\\
			\vdots&\vdots&\ddots&\vdots\\
			x_2^{n-1}-x_1x_2^{n-2}&x_3^{n-1}-x_1x_3^{n-2}&\cdots&x_n^{n-1}-x_1x_n^{n-2}
		\end{vmatrix}\\
		&=(x_2-x_1)\cdots(x_n-x_1)
		\begin{vmatrix}
			1&1&\cdots&1\\
			x_2&x_3&\cdots&x_n\\
			\vdots&\vdots&\ddots&\vdots\\
			x_2^{n-2}&x_3^{n-2}&\cdots&x_n^{n-2}
		\end{vmatrix}\\
		&=(x_2-x_1)\cdots(x_n-x_1)V(x_2,\cdots,x_n)\\
		&=\prod_{i<j}(x_j-x_i).
	\end{aligned}
	\]
\end{example}


\begin{corollary}\label{3.8.4}
	设 $A=(a_{ij})\in M_n(\F)$，任取 $A$ 的第 $i$ 行，则对任意 $k_1,\cdots,k_n$ 有
	\[
	k_1C_{i1}+\cdots+k_nC_{in}
	=
	\begin{vmatrix}
		a_{11}&\cdots&a_{1n}\\
		\vdots&&\vdots\\
		a_{i-1,1}&\cdots&a_{i-1,n}\\
		k_1&\cdots&k_n\\
		a_{i+1,1}&\cdots&a_{i+1,n}\\
		\vdots&&\vdots\\
		a_{n1}&\cdots&a_{nn}
	\end{vmatrix}.
	\]
	对列形式同理。
\end{corollary}

\begin{corollary}\label{3.8.5}
	在 $|A|$ 中，$n>2$，取第 $i,j$ 行，有
	\[
	\bigl(A(i,1)\ \cdots\ A(i,n)\bigr)
	\begin{pmatrix}
		C_{j1}\\
		\vdots\\
		C_{jn}
	\end{pmatrix}
	=\delta_{ij}|A|.
	\]
\end{corollary}

\begin{definition}\label{3.8.6}\phantomsection\label{p1:concept:c3-l2481}
	对 $A\in M_n(\F)$，其伴随矩阵 $A^*$ 定义为
	\[
	A^*(i,j)=C_{ji}.
	\]
	也即
	\[
	A=\begin{pmatrix}
		a_{11}&\cdots&a_{1n}\\
		\vdots&\ddots&\vdots\\
		a_{n1}&\cdots&a_{nn}
	\end{pmatrix}
	\quad\Longrightarrow\quad
	A^*=\begin{pmatrix}
		C_{11}&\cdots&C_{n1}\\
		\vdots&\ddots&\vdots\\
		C_{1n}&\cdots&C_{nn}
	\end{pmatrix}.
	\]
\end{definition}

\begin{theorem}[Laplace 定理的矩阵形式]\label{3.8.7}
	\[
	A^*A=AA^*=|A|E_n.
	\]
\end{theorem}

\begin{proof}
	由 \ref{3.8.5}，取 $i=j$ 即得 $AA^*=|A|E_n$，用 \ref{3.8.5} 相应的列形式，即得 $A^*A=|A|E_n$。
	
	
\end{proof}
当 $A$ 可逆时，
\[
A^{-1}=A^*|A|^{-1}.
\]
考虑 $|A^*|$，由 $A^*A=|A|E_n$，得
\[
|A^*||A|=|A|^n,
\qquad |A^*|=|A|^{n-1}.
\]
\begin{proposition}\label{3.8.8}
	设 $A\in M_n(\F)$，则
	\[
	r(A^*)=
	\begin{cases}
		n, & r(A)=n,\\
		1, & r(A)=n-1,\\
		0, & r(A)<n-1.
	\end{cases}
	\]
\end{proposition}

\begin{proof}
	当 $r(A)=n$ 时，$r(A^{-1})=n$，由 $A^{-1}=A^*|A|^{-1}$，故 $r(A^*)=n$。
	当 $r(A)<n$ 时，由 $AA^*=0$，从而由 \ref{3.5.5} 有
	\[
	r(A)+r(A^*)\leq n.
	\]
	若 $r(A)=n-1$，则 $r(A^*)\leq1$，而 $A^*$ 不是零矩阵，故 $r(A^*)=1$。
	若 $r(A)<n-1$，则 $A^*=0$，故 $r(A^*)=0$。
\end{proof}

\begin{definition}\label{3.8.9}\phantomsection\label{p1:concept:c3-l2544}
	取
	\[
	i_1<\cdots<i_k,
	\qquad
	j_1<\cdots<j_k.
	\]
	用
	\[
	A(i_1,\cdots,i_k\mid j_1,\cdots,j_k)
	\]
	表示取出 $A$ 的 $(i_p,j_q)$ 元构成的矩阵，称为 $A$ 的子矩阵。
\end{definition}

\begin{theorem}(Laplace 定理的子矩阵形式)\label{3.8.10}
	设
	\[
	i_1<\cdots<i_k,
	\]
	则
	\[
	\begin{aligned}
		|A|
		&=(-1)^{i_1+\cdots+i_k}
		\sum_{1\leq j_1<\cdots<j_k\leq n}
		(-1)^{j_1+\cdots+j_k}
		\left|A(i_1,\cdots,i_k\mid j_1,\cdots,j_k)\right|\\
		&\qquad\qquad\qquad\cdot
		\left|A(i'_1,\cdots,i'_{n-k}\mid j'_1,\cdots,j'_{n-k})\right|.
	\end{aligned}
	\]
	其中
	\[
	1\leq i'_1<\cdots<i'_{n-k}\in\{1,2,\ldots,n\}\setminus\{i_1,\cdots,i_k\}
	\]
	由 $i_1,\cdots,i_k$ 唯一确定，
	\[
	1\leq j'_1<\cdots<j'_{n-k}\in\{1,2,\ldots,n\}\setminus\{j_1,\cdots,j_k\}
	\]
	由 $j_1,\cdots,j_k$ 唯一确定。
\end{theorem}

下面，考虑分块矩阵的行列式。

\begin{proposition}\label{3.8.11}
	若 $A,C$ 是方阵，则
	\[
	\begin{vmatrix}
		A&B\\
		0&C
	\end{vmatrix}
	=|A|\,|C|.
	\]
\end{proposition}

\begin{proof}
	当 $A$ 可逆时，
	\[
	\begin{pmatrix}
		A&0\\
		0&E
	\end{pmatrix}
	\begin{pmatrix}
		E&A^{-1}B\\
		0&C
	\end{pmatrix}
	=
	\begin{pmatrix}
		A&B\\
		0&C
	\end{pmatrix},
	\]
	即得结论。一般情形可由扰动法得到。
\end{proof}

更进一步，对准上三角阵，
\[
\begin{vmatrix}
	A_{11}&*&\cdots&*\\
	0&A_{22}&\cdots&*\\
	\vdots&\vdots&\ddots&\vdots\\
	0&0&\cdots&A_{nn}
\end{vmatrix}
=|A_{11}|\,|A_{22}|\cdots|A_{nn}|.
\]

\section{一些应用与习题}

\begin{definition}\label{3.9.1}\phantomsection\label{p1:concept:c3-l2632}
	设 $A\in M_{m\times n}(\F)$。$A$ 的任意 $s\times s$ 型子矩阵的行列式称为 $A$ 的一个子式。
	若取行指标 $i_1<\cdots<i_s$，列指标 $j_1<\cdots<j_s$，则相应的 $s$ 阶子式记为
	\[
	\left|A(i_1,\cdots,i_s\mid j_1,\cdots,j_s)\right|.
	\]
\end{definition}

\begin{proposition}\label{3.9.2}
	设 $A\in M_{m\times n}(\F)$，则
	\begin{enumerate}
		\item $r(A)\geq r\Longleftrightarrow A$ 至少有一个 $r$ 阶子式不为 $0$；
		\item $r(A)\leq r\Longleftrightarrow A$ 所有阶数大于 $r$ 的子式全为 $0$。
	\end{enumerate}
	特别地，
	\[
	r(A)=r
	\Longleftrightarrow
	A\text{ 至少有一个 }r\text{ 阶子式不为 }0,
	\text{ 且所有大于 }r\text{ 阶的子式全为 }0.
	\]
\end{proposition}

\begin{proof}
	$1.$ 若 $r(A)\geq r$，则 $A$ 有 $r$ 个线性无关的行
	\[
	\beta_{j_1},\cdots,\beta_{j_r}.
	\]
	取出子矩阵
	\[
	A_1=\begin{pmatrix}
		\beta_{j_1}\\
		\vdots\\
		\beta_{j_r}
	\end{pmatrix},
	\]
	那么 $r(A_1)\geq r$，故 $A_1$ 有 $r$ 个无关的列。设为第 $k_1,\cdots,k_r$ 列，取出
	\[
	A(j_1,\cdots,j_r\mid k_1,\cdots,k_r)\in M_r(\F).
	\]
	这个矩阵满秩，且是一个方阵，即 $A$ 有一个 $r$ 阶子式不为 $0$。
	当 $r(A)=r$ 时，$A$ 也至少有一个 $r$ 阶子式不为 $0$。
	
	$2.$ 若 $r(A)\leq r$，先考虑 $r(A)<r$。设 $A$ 有一个 $r$ 阶子式不为 $0$，即
	\[
	|A_2|=\left|A(j_1,\cdots,j_r\mid k_1,\cdots,k_r)\right|\neq0.
	\]
	则 $A_2$ 的行向量无关。对应地，这些行向量是 $A$ 行向量的缩短组，那么 $A$ 的相应行向量无关，即 $r(A)\geq r$，矛盾。
	故 $r(A)<r$ 时，任意 $r$ 阶子式为 $0$，从而任意大于 $r$ 阶的子式为 $0$。
	
	对 $r(A)=r$，取
	\[
	A(j_1,\cdots,j_{r+1}\mid k_1,\cdots,k_{r+1})
	\]
	的行向量，必然线性相关，其子式为 $0$。从而大于 $r$ 阶的子式为 $0$。
\end{proof}

在这一命题中，可以再定义 $A$ 的秩 $r(A)$。在一些情况下时有用的：
\[
\text{$A$ 的秩 }r(A)\text{ 定义为 $A$ 的非零子式的最高阶数。}
\]

\begin{proposition}(Cramer 法则)\label{3.9.3}
	设 $A\in M_n(\F)$ 可逆，则 $Ax=\beta$ 的唯一解为
	\[
	x=\frac{1}{|A|}
	\begin{pmatrix}
		D_1\\
		\vdots\\
		D_n
	\end{pmatrix},
	\]
	其中 $D_j$ 是用 $\beta$ 代替 $A$ 的第 $j$ 列后得到的 $n$ 阶行列式。
\end{proposition}

\begin{proof}
	由 $A$ 可逆，$Ax=\beta$ 的唯一解为
	\[
	x=A^{-1}\beta
	=\frac{1}{|A|}A^*\beta.
	\]
	设 $\beta=(b_1,\cdots,b_n)^T$，则
	\[
	\begin{aligned}
		x
		&=\frac{1}{|A|}
		\begin{pmatrix}
			C_{11}&C_{21}&\cdots&C_{n1}\\
			C_{12}&C_{22}&\cdots&C_{n2}\\
			\vdots&\vdots&\ddots&\vdots\\
			C_{1n}&C_{2n}&\cdots&C_{nn}
		\end{pmatrix}
		\begin{pmatrix}
			b_1\\
			\vdots\\
			b_n
		\end{pmatrix}\\
		&=\frac{1}{|A|}
		\begin{pmatrix}
			C_{11}b_1+\cdots+C_{n1}b_n\\
			\vdots\\
			C_{1n}b_1+\cdots+C_{nn}b_n
		\end{pmatrix}.
	\end{aligned}
	\]
	由推论 \ref{3.8.4}，上式每一个分量正是用 $\beta$ 替换 $A$ 的相应列后所得的行列式 $D_j$，故
	\[
	x=\frac{1}{|A|}
	\begin{pmatrix}
		D_1\\
		\vdots\\
		D_n
	\end{pmatrix}.
	\]
\end{proof}

\begin{problem}
 若一个方阵的每个元素都介于 $0$ 与 $1$ 之间，并且每一行元素之和均为 $1$，则称该方阵为 \textbf{Markov 矩阵}。证明：Markov 矩阵的乘积仍为 Markov 矩阵。
\end{problem}

\begin{proof}
	设 $A=(a_{ij})$、$B=(b_{ij})$ 是 $n$ 阶 Markov 矩阵，并令
	\[
	C=(c_{ij})=AB.
	\]
	固定 $j$，记
	\[
	b_j=\max\{b_{kj}:1\leq k\leq n\}\leq1.
	\]
	由于 $a_{ik},b_{kj}\geq0$ 且 $\sum_{k=1}^{n}a_{ik}=1$，对任意 $1\leq i,j\leq n$，有
	\begin{equation}\tag{1}
		0\leq c_{ij}=\sum_{k=1}^{n}a_{ik}b_{kj}
		\leq b_j\sum_{k=1}^{n}a_{ik}=b_j\leq1.
	\end{equation}
	另一方面，对任意 $1\leq i\leq n$，交换有限求和的次序可得
	\begin{equation}\tag{2}
		\begin{aligned}
			\sum_{j=1}^{n}c_{ij}
			&=\sum_{j=1}^{n}\sum_{k=1}^{n}a_{ik}b_{kj}
			=\sum_{k=1}^{n}a_{ik}\left(\sum_{j=1}^{n}b_{kj}\right)\\
			&=\sum_{k=1}^{n}a_{ik}=1.
		\end{aligned}
	\end{equation}
	由 $(1),(2)$，$C=AB$ 是 Markov 矩阵。
\end{proof}

\begin{problem}
 设 $A,B$ 为实方阵，且
	\[
	A^{T}=A,\qquad B^{T}=-B,\qquad A^{2}=B^{2}.
	\]
	证明 $A=B=0$。
\end{problem}

\begin{proof}
	令 $C=A^{2}=B^{2}$。对任意 $1\leq i\leq n$，由 $A(i,j)=A(j,i)$，得
	\begin{equation}\tag{1}
		C(i,i)=A^{2}(i,i)
		=\sum_{j=1}^{n}A(i,j)A(j,i)
		=\sum_{j=1}^{n}A(i,j)^{2}.
	\end{equation}
	由 $B(j,i)=-B(i,j)$，得
	\begin{equation}\tag{2}
		C(i,i)=B^{2}(i,i)
		=\sum_{j=1}^{n}B(i,j)B(j,i)
		=-\sum_{j=1}^{n}B(i,j)^{2}.
	\end{equation}
	由 $(1),(2)$，
	\begin{equation}\tag{3}
		\sum_{j=1}^{n}A(i,j)^{2}+\sum_{j=1}^{n}B(i,j)^{2}=0.
	\end{equation}
	各项均为非负实数，故 $A(i,j)=B(i,j)=0$。于是 $A=B=0$。
\end{proof}

\begin{problem}
	\begin{enumerate}
		\item 设 $f,g:V_1\to V_2$ 为线性映射，并定义
		\[
		\operatorname{im}(f)+\operatorname{im}(g)
		=\{\gamma:\gamma=\alpha+\beta,
		\ \alpha\in\operatorname{im}(f),\ \beta\in\operatorname{im}(g)\}.
		\]
		证明
		\[
		\operatorname{im}(f+g)\subseteq
		\operatorname{im}(f)+\operatorname{im}(g),
		\qquad
		\ker f\cap\ker g\subseteq\ker(f+g).
		\]
		
		\item 设 $f:V_1\to V_2$、$g:V_2\to V_3$ 为线性映射。若 $g\circ f=0$，证明
		\[
		\operatorname{im}(f)\subseteq\ker g.
		\]
	\end{enumerate}
\end{problem}

\begin{proof}
	$(1)$ 任取 $\gamma\in\operatorname{im}(f+g)$，则存在 $\alpha\in V_1$，使得
	\[
	\gamma=(f+g)(\alpha)=f(\alpha)+g(\alpha)
	\in\operatorname{im}(f)+\operatorname{im}(g).
	\]
	故
	\[
	\operatorname{im}(f+g)\subseteq
	\operatorname{im}(f)+\operatorname{im}(g).
	\]
	任取 $\xi\in\ker f\cap\ker g$，则 $f(\xi)=g(\xi)=0$，从而
	\[
	(f+g)(\xi)=f(\xi)+g(\xi)=0,
	\]
	故 $\xi\in\ker(f+g)$，即
	\[
	\ker f\cap\ker g\subseteq\ker(f+g).
	\]
	
	$(2)$ 任取 $\alpha\in\operatorname{im}(f)$，则存在 $\beta\in V_1$，使得 $\alpha=f(\beta)$。由 $g\circ f=0$，
	\[
	g(\alpha)=g(f(\beta))=(g\circ f)(\beta)=0.
	\]
	因此 $\alpha\in\ker g$，从而
	\[
	\operatorname{im}(f)\subseteq\ker g.
	\]
\end{proof}

\begin{problem}
 设 $A$ 为 $m\times n$ 矩阵，$B$ 为 $n\times\ell$ 矩阵。证明
	\[
	r(AB)\geq r(A)+r(B)-n.
	\]
\end{problem}

\begin{proof}
	考虑分块矩阵
	\[
	X=\begin{pmatrix}
		E_n&B\\
		A&0
	\end{pmatrix}.
	\]
	$A$ 与 $B$ 分别位于互不相交的行、列块中，故
	\begin{equation}\tag{1}
		r(X)\geq r(A)+r(B).
	\end{equation}
	对 $X$ 作可逆的分块初等变换：
	\[
	\begin{aligned}
		Y
		&=\begin{pmatrix}
			E_n&0\\
			-A&E_m
		\end{pmatrix}
		\begin{pmatrix}
			E_n&B\\
			A&0
		\end{pmatrix}
		\begin{pmatrix}
			E_n&-B\\
			0&E_{\ell}
		\end{pmatrix}\\
		&=\begin{pmatrix}
			E_n&0\\
			0&-AB
		\end{pmatrix}.
	\end{aligned}
	\]
	左右两个矩阵均可逆，故 $r(Y)=r(X)$。于是
	\[
	n+r(AB)=r(Y)=r(X)\geq r(A)+r(B),
	\]
	即
	\[
	r(AB)\geq r(A)+r(B)-n.
	\]
\end{proof}

\begin{problem}
设
	\[
	A=(a_{ij})_{m\times n},\qquad B=(b_{ij})_{s\times t}
	\]
	\phantomsection\label{p1:concept:c3-l2915}是数域 $\F$ 上的两个矩阵。定义它们的 \textbf{Kronecker 积}（或张量积）为如下 $ms\times nt$ 分块矩阵：
	\[
	A\otimes B=
	\begin{pmatrix}
		a_{11}B&\cdots&a_{1n}B\\
		\vdots&&\vdots\\
		a_{m1}B&\cdots&a_{mn}B
	\end{pmatrix}.
	\]
	证明：
	\begin{enumerate}
		\item 对任意 $k\in\F$，
		\[
		(kA)\otimes B=k(A\otimes B),
		\qquad
		A\otimes(kB)=k(A\otimes B).
		\]
		
		\item
		\[
		(A\otimes B)^T=A^T\otimes B^T,
		\]
		但一般并没有 $A\otimes B=B\otimes A$。
		
		\item 当矩阵乘法有意义时，
		\[
		(A\otimes B)(C\otimes D)=AC\otimes BD.
		\]
		
		\item 若 $A,B$ 均可逆，则 $A\otimes B$ 可逆，且
		\[
		(A\otimes B)^{-1}=A^{-1}\otimes B^{-1}.
		\]
		
		\item 若 $A$ 为 $n\times n$ 矩阵，$B$ 为 $m\times m$ 矩阵，则
		\[
		\det(A\otimes B)=(\det A)^m(\det B)^n.
		\]
	\end{enumerate}
\end{problem}

\begin{proof}
	$(1)$ 由分块矩阵的定义直接得到。
	
	$(2)$ 将 $A\otimes B$ 的第 $(i,j)$ 个块 $a_{ij}B$ 转置后，位于第 $(j,i)$ 个块，且
	\[
	(a_{ij}B)^T=a_{ij}B^T,
	\]
	故 $(A\otimes B)^T=A^T\otimes B^T$。一般交换律不成立。例如
	\[
	A=\begin{pmatrix}1&0\\0&0\end{pmatrix},
	\qquad
	B=\begin{pmatrix}0&1\\0&0\end{pmatrix},
	\]
	则 $A\otimes B$ 的非零元位于第 $(1,2)$ 个位置，而 $B\otimes A$ 的非零元位于第 $(1,3)$ 个位置，故二者不相等。
	
	$(3)$ 设 $A=(a_{ij})$、$C=(c_{jk})$。乘积的第 $(i,k)$ 个块为
	\[
	\sum_j(a_{ij}B)(c_{jk}D)
	=\left(\sum_j a_{ij}c_{jk}\right)BD,
	\]
	这正是 $AC\otimes BD$ 的第 $(i,k)$ 个块。
	
	$(4)$ 由 $(3)$，
	\[
	(A\otimes B)(A^{-1}\otimes B^{-1})
	=E_n\otimes E_m=E_{nm},
	\]
	反向乘积同理，故结论成立。
	
	$(5)$ 由 $(3)$，
	\[
	A\otimes B=(A\otimes E_m)(E_n\otimes B).
	\]
	其中 $E_n\otimes B$ 是含有 $n$ 个对角块 $B$ 的分块对角阵，故
	\[
	\det(E_n\otimes B)=(\det B)^n.
	\]
	又存在一个置换矩阵 $P$，使得
	\[
	P^{-1}(A\otimes E_m)P=E_m\otimes A.
	\]
	因此
	\[
	\det(A\otimes E_m)=\det(E_m\otimes A)=(\det A)^m.
	\]
	两式相乘即得
	\[
	\det(A\otimes B)=(\det A)^m(\det B)^n.
	\]
\end{proof}

\begin{problem}
设 $A,B$ 为 $n\times n$ 矩阵。若 $E_n+AB$ 可逆，证明 $E_n+BA$ 也可逆，并且
	\[
	(E_n+BA)^{-1}=E_n-B(E_n+AB)^{-1}A.
	\]
\end{problem}

\begin{proof}
	记
	\[
	X=E_n-B(E_n+AB)^{-1}A.
	\]
	则
	\[
	\begin{aligned}
		(E_n+BA)X
		&=E_n-B(E_n+AB)^{-1}A+BA
		-BAB(E_n+AB)^{-1}A\\
		&=E_n+BA
		-B\bigl((E_n+AB)^{-1}+AB(E_n+AB)^{-1}\bigr)A\\
		&=E_n+BA
		-B\bigl((E_n+AB)(E_n+AB)^{-1}\bigr)A\\
		&=E_n+BA-BE_nA\\
		&=E_n.
	\end{aligned}
	\]
	因此 $X$ 是方阵 $E_n+BA$ 的左逆，故 $E_n+BA$ 可逆，且
	\[
	(E_n+BA)^{-1}=X
	=E_n-B(E_n+AB)^{-1}A.
	\]
\end{proof}


形式上，考虑
\[\begin{aligned}
	(1+ba)^{-1}&=1-(ba)+(ba)^2-(ba)^3+\cdots\\
	&=1-ba+b(ab)a-b(ab)^2a+\cdots\\
	&=1-b(1-ab+(ab)^2+\cdots)a\\
	&=1-b(1+ab)^{-1}a
\end{aligned}  \]
可以猜到上面的公式

\begin{problem}
设 $A,B$ 为方阵，且 $A,B,AB-E_n$ 均可逆。证明 $A-B^{-1}$ 与
	\[
	(A-B^{-1})^{-1}-A^{-1}
	\]
	均可逆，并且
	\[
	\left((A-B^{-1})^{-1}-A^{-1}\right)^{-1}=ABA-A.
	\]
\end{problem}

\begin{proof}
	由
	\[
	(A-B^{-1})B=AB-E_n
	\]
	以及 $B$、$AB-E_n$ 均可逆，知 $A-B^{-1}$ 可逆。对上式两边取逆，得
	\[
	B^{-1}(A-B^{-1})^{-1}=(AB-E_n)^{-1},
	\]
	从而
	\[
	(A-B^{-1})^{-1}=B(AB-E_n)^{-1}.
	\]
	于是
	\[
	\begin{aligned}
		&\left((A-B^{-1})^{-1}-A^{-1}\right)(ABA-A)\\
		&\quad=\left(B(AB-E_n)^{-1}-A^{-1}\right)(AB-E_n)A\\
		&\quad=B(AB-E_n)^{-1}(AB-E_n)A
		-A^{-1}(AB-E_n)A\\
		&\quad=BA-(B-A^{-1})A\\
		&\quad=E_n.
	\end{aligned}
	\]
	故 $(A-B^{-1})^{-1}-A^{-1}$ 可逆，且其逆矩阵为 $ABA-A$。
\end{proof}

\begin{problem}
设
	\[
	P_1=(a_1,a_2),\qquad P_2=(b_1,b_2)
	\]
	是实平面 $\R^2$ 中与原点 $O$ 不共线的两点。设 $\varphi$ 是由 $P_1,P_2,O$ 确定的平行四边形。证明 $\varphi$ 的面积等于行列式
	\[
	\begin{vmatrix}
		a_1&a_2\\
		b_1&b_2
	\end{vmatrix}
	\]
	的绝对值，并将此结论推广到 $\R^n\ (n>2)$。
\end{problem}

\begin{proof}
	设 $\alpha$ 为 $\overrightarrow{OP_1}$ 与 $\overrightarrow{OP_2}$ 的夹角。由余弦定理与内积公式，平行四边形的面积 $S$ 为
	\[
	\begin{aligned}
		S
		&=\lVert\overrightarrow{OP_1}\rVert
		\lVert\overrightarrow{OP_2}\rVert\sin\alpha\\
		&=\sqrt{
			\lVert\overrightarrow{OP_1}\rVert^2
			\lVert\overrightarrow{OP_2}\rVert^2
			-(\overrightarrow{OP_1}\cdot\overrightarrow{OP_2})^2}\\
		&=\sqrt{(a_1^2+a_2^2)(b_1^2+b_2^2)
			-(a_1b_1+a_2b_2)^2}\\
		&=\sqrt{(a_1b_2-a_2b_1)^2}\\
		&=|a_1b_2-a_2b_1|\\
		&=\left|
		\begin{vmatrix}
			a_1&a_2\\
			b_1&b_2
		\end{vmatrix}
		\right|.
	\end{aligned}
	\]
	
	一般地，设 $\alpha_i=(a_{i1},\ldots,a_{in})\in\R^n$，$1\leq i\leq n$，并令 $M=(a_{ij})$。由 Gram--Schmidt 正交化，$\alpha_1,\ldots,\alpha_n$ 张成的平行多面体体积满足
	\[
	V^2=\det(MM^T)=(\det M)^2.
	\]
	因此
	\[
	V=|\det M|.
	\]
	当 $n=3$ 时，这就是由三个向量张成的平行六面体体积公式。
\end{proof}

\begin{problem}
设 $A$ 为奇数阶方阵。证明
	\[
	\det(A-A^T)=0.
	\]
\end{problem}

\begin{proof}
	令 $X=A-A^T$，则 $X^T=-X$，即 $X$ 为反对称矩阵。若 $X$ 为 $n$ 阶且 $n$ 为奇数，则
	\[
	\det X=\det X^T=\det(-X)=(-1)^n\det X=-\det X.
	\]
	故 $\det X=0$，从而
	\[
	\det(A-A^T)=0.
	\]
\end{proof}

\begin{problem}
令
	\[
	s_k=x_1^k+x_2^k+\cdots+x_n^k,
	\qquad x_i\in\Z,\quad k\geq0,
	\]
	并设
	\[
	A=\begin{pmatrix}
		s_0&s_1&s_2&\cdots&s_{n-1}\\
		s_1&s_2&s_3&\cdots&s_n\\
		s_2&s_3&s_4&\cdots&s_{n+1}\\
		\vdots&\vdots&\vdots&\ddots&\vdots\\
		s_{n-1}&s_n&s_{n+1}&\cdots&s_{2n-2}
	\end{pmatrix}.
	\]
	证明
	\[
	\det A=\prod_{1\leq i<j\leq n}(x_i-x_j)^2.
	\]
\end{problem}

\begin{proof}
	考虑 Vandermonde 矩阵
	\[
	V(x_1,x_2,\ldots,x_n)
	=\begin{pmatrix}
		1&1&\cdots&1\\
		x_1&x_2&\cdots&x_n\\
		\vdots&\vdots&\ddots&\vdots\\
		x_1^{n-1}&x_2^{n-1}&\cdots&x_n^{n-1}
	\end{pmatrix}.
	\]
	由 $s_k=\sum_{j=1}^{n}x_j^k$，逐项计算可得
	\[
	A=V(x_1,\ldots,x_n)V(x_1,\ldots,x_n)^T.
	\]
	因此
	\[
	\begin{aligned}
		\det A
		&=\det\bigl(VV^T\bigr)\\
		&=(\det V)^2\\
		&=\prod_{1\leq i<j\leq n}(x_i-x_j)^2.
	\end{aligned}
	\]
\end{proof}

\begin{problem}
设可逆矩阵 $A$ 的每一行元素之和均为 $k$。证明：$A^{-1}$ 的每一行元素之和均为 $1/k$。
\end{problem}

\begin{proof}
	令
	\[
	\delta=(1,1,\ldots,1)^T\in\F^n.
	\]
	对任意矩阵 $B=(b_{ij})_{n\times n}$，向量 $B\delta$ 的第 $i$ 个分量正是 $B$ 的第 $i$ 行元素之和。因此，$A$ 的各行元素之和均为 $k$ 等价于
	\[
	A\delta=k\delta.
	\]
	因为 $A$ 可逆且 $\delta\neq0$，必有 $k\neq0$。在等式两边左乘 $A^{-1}$，得
	\[
	\delta=kA^{-1}\delta,
	\qquad\text{即}
	\qquad
	A^{-1}\delta=\frac1k\delta.
	\]
	这表明 $A^{-1}$ 的每一行元素之和均为 $1/k$。
\end{proof}

\begin{problem}
	\textbf{分块行列式公式.}\quad 设 $A,B,C,D$ 均为同阶方阵。若 $A$ 可逆且 $AC=CA$，证明
	\[
	\begin{vmatrix}
		A&B\\
		C&D
	\end{vmatrix}
	=\det(AD-CB).
	\]
\end{problem}

\begin{proof}
	对分块矩阵作可逆的分块初等行变换：
	\[
	\begin{pmatrix}
		E_n&0\\
		-CA^{-1}&E_n
	\end{pmatrix}
	\begin{pmatrix}
		A&B\\
		C&D
	\end{pmatrix}
	=
	\begin{pmatrix}
		A&B\\
		0&D-CA^{-1}B
	\end{pmatrix}.
	\]
	左侧第一个矩阵的行列式为 $1$，故
	\begin{equation}\tag{1}
		\begin{vmatrix}
			A&B\\
			C&D
		\end{vmatrix}
		=\det A\,\det(D-CA^{-1}B).
	\end{equation}
	由 $AC=CA$ 及 $A$ 可逆，得 $ACA^{-1}=C$。于是
	\begin{equation}\tag{2}
		\begin{aligned}
			\det A\,\det(D-CA^{-1}B)
			&=\det\bigl(A(D-CA^{-1}B)\bigr)\\
			&=\det(AD-ACA^{-1}B)\\
			&=\det(AD-CB).
		\end{aligned}
	\end{equation}
	由 $(1),(2)$ 即得结论。
\end{proof}

\begin{remark}
	若在相应的交换条件下需要处理不可逆情形，可考虑以 $A+tE_n$ 代替 $A$，先对除有限个 $t$ 之外的情形应用上述公式，再利用行列式关于 $t$ 的多项式恒等性令 $t\to0$。
\end{remark}

\begin{problem}
 设 $n,m\geq1$，且
	\[
	f(x)=a_0x^n+a_1x^{n-1}+\cdots+a_n,
	\qquad
	g(x)=b_0x^m+b_1x^{m-1}+\cdots+b_m,
	\]
	\phantomsection\label{p1:concept:c3-l3286}其中 $a_0b_0\neq0$。定义 $f,g$ 的结式矩阵为如下 $(m+n)\times(m+n)$ 矩阵：
	\[
	R(f,g)=
	\begin{pmatrix}
		a_0&a_1&\cdots&a_n&0&\cdots&0\\
		0&a_0&a_1&\cdots&a_n&\ddots&\vdots\\
		\vdots&\ddots&\ddots&\ddots&&\ddots&0\\
		0&\cdots&0&a_0&a_1&\cdots&a_n\\[2mm]
		b_0&b_1&\cdots&b_m&0&\cdots&0\\
		0&b_0&b_1&\cdots&b_m&\ddots&\vdots\\
		\vdots&\ddots&\ddots&\ddots&&\ddots&0\\
		0&\cdots&0&b_0&b_1&\cdots&b_m
	\end{pmatrix},
	\]
	其中上面由 $f$ 的系数组成的移位行共有 $m$ 行，下面由 $g$ 的系数组成的移位行共有 $n$ 行。证明：$f,g$ 互素，当且仅当 $R(f,g)$ 可逆。
\end{problem}

\begin{proof}
	记 $d=(f,g)$。若 $f,g$ 不互素，则 $\deg d>0$。令
	\[
	f_1(x)=\frac{f(x)}{d(x)},
	\qquad
	g_1(x)=\frac{g(x)}{d(x)}.
	\]
	则
	\[
	f_1(x)g(x)=g_1(x)f(x),
	\]
	且
	\[
	\deg f_1<n,
	\qquad
	\deg g_1<m.
	\]
	把 $f_1,g_1$ 的系数依次排成列向量，比较上式两边各次幂的系数，即得到一个以 $R(f,g)^T$ 为系数矩阵的非零齐次解。因此 $R(f,g)$ 不可逆。
	
	反之，若 $R(f,g)$ 不可逆，则存在不全为零的多项式 $u,v$，满足
	\[
	\deg u<m,
	\qquad
	\deg v<n,
	\qquad
	u(x)f(x)+v(x)g(x)=0.
	\]
	这里的等式正是齐次方程组
	\[
	R(f,g)^T
	\begin{pmatrix}
		\text{$u$ 的系数}\\
		\text{$v$ 的系数}
	\end{pmatrix}=0
	\]
	所表达的系数关系。若 $f,g$ 互素，则由 $f\mid vg$ 可得 $f\mid v$。但 $\deg v<n=\deg f$，故 $v=0$；再由 $uf=0$ 得 $u=0$，与二者不全为零矛盾。因此 $f,g$ 不互素。
	
	综上，$f,g$ 互素当且仅当 $R(f,g)$ 可逆。
\end{proof}

